%\documentclass[manuscript,blind]{geophysics}
%\documentclass[paper,twocolumn,twoside]{geophysics}
\documentclass[manuscript,noblind]{geophysics}

% An example of defining macros
\newcommand{\rs}[1]{\mathstrut\mbox{\scriptsize\rm #1}}
\newcommand{\rr}[1]{\mbox{\rm #1}}

%% Extra packages
\usepackage{amsmath}
\usepackage{amssymb}
\usepackage{amsthm}
\usepackage{nicefrac}
\usepackage{array}
\usepackage[linesnumbered,lined,ruled,commentsnumbered]{algorithm2e}

% replace the term "Input" by "Initialization" in algorithms
\SetKwInOut{Input}{Initialization}

%% For creating theorems, corollaries, lemmas and definitions
%% source: https://www.overleaf.com/learn/latex/Theorems_and_proofs
%\newtheorem{theorem}{Theorem}[section]
\newtheorem{theorem}{Theorem}
%\newtheorem{corollary}{Corollary}[theorem]
%\newtheorem{proposition}{Proposition}[theorem]
\newtheorem{proposition}{Proposition}
%\newtheorem{lemma}{Lemma}[theorem]
%\newtheorem{lemma}{Lemma}
%\theoremstyle{definition}%
\newtheorem{definition}{Definition}

%% Define the QED symbol as a black square
\renewcommand\qedsymbol{$\blacksquare$}

% for creating tables
%\usepackage{booktabs}
%\usepackage{array}

\begin{document}
	\title{Potential-field theoretical insights on planar equivalent layers}
	
	\renewcommand{\thefootnote}{\fnsymbol{footnote}} 
	
	\ms{GEO-2024XXXX} % manuscript number
	
	\address{
		\footnotemark[1] Corresponding author: vanderlei@on.br
		\footnotemark[2] Observat\'orio Nacional, Rio de Janeiro, Brazil\\
		\footnotemark[3] Universidade do Estado do Rio de Janeiro, Rio de Janeiro, Brazil\\
		\footnotemark[4] TU Bergakademie Freiberg, Institute of Geophysics and Geoinformatics, Freiberg, Germany
	}
	\author{Vanderlei C. Oliveira Jr.\footnotemark[2]\footnotemark[1], Christian Gerhards\footnotemark[4]
		Val\'eria C. F. Barbosa\footnotemark[2], Andr\'e L. A. Reis \footnotemark[3],
		Edson F. Luza\footnotemark[2], Xinpeng Huang\footnotemark[4]}
	
	%\footer{Example}
	\lefthead{Oliveira Jr. et al.}
	\righthead{Planar equivalent layers}
	
	\maketitle
	
	% Abstract
	\begin{abstract}
		We present a comprehensive theoretical framework for potential-field transformations 
		using planar equivalent layers, a technique empirically developed since the 1960s in applied
		geophysics. We aim to replace the widely accepted equivalent-layer principle – primarily
		based on empirical generalizations – with a more rigorous theoretical foundation. We 
		introduce and mathematically prove the \textit{planar equivalent-layer theorem}, which establishes
		the existence of continuous planar layers that can accurately reproduce potential-field data
		within a predefined error level, even for layers that cross the true sources.
		We present closed-form expressions for planar equivalent layers by modifying the solutions of
		Dirichlet’s and Neumann’s problems and demonstrate that planar equivalent layers are not merely
		mathematical constructs, but have tangible physical significance related to boundary value problems 
		in potential theory.
		Using our closed-form expressions, we mathematically show the equivalence between potential-field
		transformations made via equivalent-layer technique and those commonly derived from the Fourier
		transform of observed potential-field data.
		Moreover, we define different classes of equivalent layers and their relationships with 
		the observed potential-field data, true sources and the desired potential-field transformation.
		Our work advances the understanding of why the equivalent-layer technique has been
		effectively applied in potential-field transformations.
	\end{abstract}
	
	% Introduction
	\section{Introduction}
	
	The equivalent-layer technique was introduced during the 1960s and has become a versatile
	tool for processing and interpreting potential-field data in applied geophysics. 
	It consists in approximating a discrete set of potential-field data 
	by a linear combination of harmonic functions representing the potential field 
	due to a discrete set of virtual sources (e.g., monopoles, dipoles, prisms).
	These virtual sources are called \textit{equivalent sources} and are commonly disposed on a surface 
	forming a discrete \textit{equivalent layer}, which is ``equivalent" 
	in the sense that it satisfactorily reproduces the observed potential-field data.
	The general approach involves estimating the constant physical property of the equivalent sources 
	(e.g., density, mass, isotropic magnetic susceptibility, magnetization intensity, magnetic moment intensity) 
	by numerically solving a constrained linear inverse problem.
	The technique can be used to accomplish various data processing steps, such as interpolation, 
	upward (or downward) continuation, gradient data processing, merging multiple data sets, 
	estimating density models from gravity data, 
	computing the amplitude of the anomalous magnetic field, 
	transforming total-field anomalies into field components, 
	as well as estimating the magnetization direction and computing the reduction to the pole
	of magnetic data. 
	Further details on the equivalent-layer technique can be found in the recent review by 
	\citet{oliveirajr_etal2023}.
	
	Several equivalent-layer methods follow the pioneer approach presented
	by \citet{dampney1969} and rely on the numerical solution of the 
	\textit{upward continuation integral} \citep[e.g.,][]{skeels1947, peters1949, henderson-zietz1949,
		dean1958, roy1962, bhattacharyya1967, henderson1970}, which is the solution of the
	\textit{Dirichlet’s problem} or the \textit{first boundary value problem of potential
		theory} on a plane located above the sources \cite[][ p. 236]{kellogg1967}. 
	From the physical point of view, this numerical solution gives rise to a predicted potential field 
	analogous to the vertical component of the gravitational attraction due to a layer of monopoles
	or the magnetic scalar potential due to a layer of vertically-magnetized dipoles.
	This approach provides the theoretical background for computing the interpolation or upward/downward 
	continuation of any potential-field data.
	However, other field transformations that have been successfully used in practice, 
	such as joint processing of gravity-gradient data 
	\cite[e.g.,][]{barnes_lumley2011} or conversion of gravity disturbance data into the 
	gravity-gradient tensor \cite[e.g.,][]{piauilino_etal2019_seg, oliveirajr_etal2023}, are 
	empirical generalizations often assumed to be valid despite the lack of mathematical proof.
	
	Most of the well-known equivalent-layer methods for magnetic data processing are 
	empirical generalizations of the \textit{Neumann's problem} or the \textit{second boundary value problem of potential 
		theory} \cite[][ p. 246]{kellogg1967}.
	These methods usually presume, without any proof, the existence of an equivalent layer of dipoles located above 
	the sources that can be used to reduce magnetic data to the pole 
	\citep[e.g.,][]{emilia1973, bhattacharyya_chan1977, silva1986, leao-silva_1989, oliveirajr_etal2013, 
		li_etal2014, oliveirajr_etal2015, ribeiro_filho_etal2020, reis_etal2020, li-etal2021, takahashi-etal2022, piauilino_etal2025}. 
	Up to now, far too little attention has been paid to the theoretical aspects related to the existence of
	a planar equivalent layer of dipoles or its relationship with a boundary value
	problem of potential theory in applied geophysics.
	By following different approaches in the wavenumber domain, 
	\citet{pedersen1991} and \citet{li_etal2014} proved the existence of a
	planar equivalent layer of dipoles with all-positive
	magnetic moment distribution that exactly reproduces the approximated total-field anomaly
	\cite[e.g., ][ p. 179]{blakely1996} caused by sources having a purely induced magnetization
	in the vertical direction.
	Their approach, however, is valid only for the case in which 
	the total-field anomaly data rely on a plane.
	\citet{reis_etal2020} generalized the works of \citet{pedersen1991} and \citet{li_etal2014}
	to prove that this planar equivalent layer exists for the case in which the dipoles have 
	the same arbitrary uniform magnetization direction of the sources, regardless of whether 
	the sources magnetization is purely induced or not.
	All those studies assume that the equivalent layer is defined in a source-free region, where  
	Laplace’s equation is obeyed. Consequently, all those studies presume that the equivalent layer 
	does not cross the true sources.
	
	The seminal work of \citet{gunn1975} was the pioneer in showing a complete description of potential-field 
	transformations in the Fourier domain. That author has also formulated such transformations in terms of 
	planar equivalent layers of monopoles and dipoles.
	However, while presuming the existence of planar layers, \citet{gunn1975} did not establish any relationship 
	with boundary-value problems in potential theory and did not take into account situations in which they 
	intersect the true sources.
	
	There have been no attempts to establish the theoretical background explaining, for example,
	the existence of a planar layer of monopoles that simultaneously fit the gravitational  potential and 
	its spatial derivatives of first and second orders; 
	the existence of a planar layer of dipoles that simultaneously fit the magnetic scalar potential and
	induction field components due to arbitrarily magnetized sources; or the existence of 
	equivalent layers that cross the sources.
	To the authors knowledge, all equivalent-layer methods found in the literature 
	consider that the layer is located above the sources; 
	a condition that cannot be ensured in practical situations with field data produced by 
	anomaly sources having undisclosed depths.
	The absence of precise depth information about the true sources, however, 
	does not prevent the real-world application of the equivalent-layer technique.
	Existing literature contains numerous examples of its use with a wide range of real-data 
	applications, from exploration to satellite geophysics \citep[e.g.,][]{gonzalez_etal2022}.
	
	Here, we present a systematic theoretical discussion on the existence of planar equivalent layers 
	that reproduce the gravitational and magnetic induction fields due to 3D sources, 
	even for the case in which the layer crosses them.
	We limit ourselves to the case in which the magnetic
	sources have a total magnetization with constant direction, but spatially-variable intensity.
	We formulate the existence of continuous planar layers as the \textit{planar equivalent-layer theorem},
	which we mathematically prove by modifying the solutions of the Dirichlet's and Neumann's problems 
	on planar surfaces that may cross the sources. 
	Our results fill the theoretical gaps in most of the existing equivalent-layer methods, including 
	those presenting alternative formulations \citep[e.g.,][]{cordell1992, guspi_novara2009}.
	We also establish how a particular equivalent layer can be used to perform a particular potential-field transformation.
	Besides, we define the theoretical link between potential-field transformations made via equivalent-layer 
	technique and those directly computed from the Fourier transform of the observed potential-field data. 
	
	We address theoretical aspects on the existence of planar equivalent layers and their relationship 
	with potential-field transformations with the purpose of replacing with the planar equivalent-layer theorem
	the widely accepted \textit{equivalent-layer principle}. 
	This principle refers to the 
	empirical generalizations by which most of the studies found in the geophysical literature justify the 
	existence of equivalent layers allowing computing potential-field transformations. 
	Part of our theoretical findings was successfully illustrated with numerical simulations 
	and an application to airborne magnetic data.
	Rather than a totally new application for the equivalent-layer technique, the main goal here is to 
	explain or shed new light on why the technique has long been successfully applied to perform potential-field transformations. 
	The present work opens the door for future investigations about, for example,  
	(i) the analysis of discretization and truncation effects applied to continuous equivalent layers; 
	(ii) the generalization of the reduction to the pole (RTP) 
	to the case in which the sources have a nonuniform total-magnetization direction;
	(iii) impact on potential-field transformations of planar non-null physical-property distributions that produce
	no field, i.e., the annihilators \citep[][ p. 66]{parker1994} or silent sources
	\citep{baratchart2013, baratchart_etal2021}; and also 
	(iv) the relationship between particular planar equivalent layers and the true 3D sources.
	
	While the focus of this work is on the theoretical foundations of the equivalent-layer technique, 
	we recognize the importance of its practical implementation. We illustrate the relevance of our results 
	through an application to field data. For a broader discussion of algorithmic choices and computational 
	strategies related to equivalent-layer methods, we refer readers to \citet{oliveirajr_etal2023}, which 
	offers a comprehensive review of these aspects. Our primary objective here is to resolve longstanding 
	theoretical gaps, particularly regarding the use of planar equivalent layers to reproduce potential-field data. 
	By establishing this foundation, we aim to deepen the understanding of why the equivalent-layer technique has 
	been so effective in geophysical practice and to support the development of future methodological advances.
	
	% Background
	\section{General background}
	\label{sec:background}
	
	\subsection{Preliminary notational and geometric definitions}
	\label{sec:preliminaries}
	
	The main symbols used throughout the present work are shown in Table \ref{tbl:Table1} and Figure \ref{fig:Figure1}.
	Some of them are briefly described in the present subsection.
	
	Let a set of 3D gravity or magnetic sources be defined over a closed region 
	$\overline{\mathcal{V}'}$ located in the subsurface, with top limited by the vertical coordinate $z^{t}_{0}$ 
	and scalar physical property distribution $\sigma(\mathbf{r}')$, $\mathbf{r}' \in \overline{\mathcal{V}'}$, 
	referred to a topocentric Cartesian coordinate system with axes 
	$x$, $y$, and $z$ pointing to north, east and down, respectively.
	A detailed definition of $\sigma(\mathbf{r}')$ is given in Table \ref{tbl:Table1}.
	
	We also consider to have potential-field observations at generic points 
	$\mathbf{r} \notin \overline{\mathcal{V}'}$ or at a finite set of points 
	$\mathbf{r}_{i} \notin \overline{\mathcal{V}'}$, $i \in \{ 1, \dots, N \}$, 
	located above the sources, i.e., with vertical coordinates satisfying $z < z^{t}_{0}$ and 
	$z_{i} < z^{t}_{0}$. These points are conveniently called \textit{observation points}.
	
	Specific horizontal surfaces are denoted by $\mathcal{S}$, with a particular superscript symbol. 
	Points and elements of area on a given horizontal surface are denoted by using the corresponding superscript symbol. 
	The constant vertical coordinate associated with a given horizontal surface is defined by using the 
	corresponding superscript symbol and subscript ``$0$''. 
	The particular surface $\mathcal{S}''$ with constant vertical coordinate $z''_{0}$ is used to define all
	planar equivalent layers presented throughout the present work. Figure \ref{fig:Figure1} shows two possible 
	definitions for $\mathcal{S}''$, one above and another crossing the sources.
	
	Finally, we use the symbol $\mathcal{F} [ \cdot ]$ to denote the two-dimensional Fourier transform of a generic scalar 
	function computed on a horizontal surface, with respect to the spatial coordinates $x$ and $y$. 
	To simplify the notation, we conveniently omit the dependence of $\mathcal{F} [ \cdot ]$ on the 
	wavenumbers $k_{x}$ and $k_{y}$.
	
	
	\subsection{Potential fields due to 3D sources}
	\label{sec:potential_fields_3D_sources}
	
	
	Let $d^{o}(\mathbf{r})$ be the potential field data produced by the 3D sources at observation points 
	$\mathbf{r}$ (Figure \ref{fig:Figure1}) as follows
	\begin{equation}
		d^{o}(\mathbf{r}) = D_{d} \, U(\mathbf{r}) \: , 
		\quad \mathbf{r} \notin \overline{\mathcal{V}'} \: ,
		\label{eq:observed-data-space}
	\end{equation}
	where $D_{d}$ is a space-domain differential operator acting on the Newtonian potential
	\begin{equation}
		U(\mathbf{r}) = \kappa \iiint\limits_{\mathcal{V}'}
		\sigma(\mathbf{r}') \: \frac{1}{\| \mathbf{r} - \mathbf{r}'\|}
		\: d\mathcal{V}' \: , 
		\quad \mathbf{r} \notin \overline{\mathcal{V}'} \: ,
		\label{eq:potential-U-space}
	\end{equation}
	$d\mathcal{V}' = dx'dy'dz'$ represents a volume element;
	$\kappa$ is a proportionality constant;
	$\tfrac{1}{\| \mathbf{r} - \mathbf{r}' \|}$ is the inverse distance function
	between the observation point $\mathbf{r}$ and a point $\mathbf{r}' \in \mathcal{V}'$
	and $\sigma(\mathbf{r}')$ is the scalar function representing 
	the density (contrast) or total-magnetization intensity within the sources (Table \ref{tbl:Table1}).
	
	The Fourier transform of the potential field data $d^{o}$ (equation \ref{eq:observed-data-space}) 
	on a horizontal surface $\mathcal{S}^{\dagger}$ above the sources (Figure \ref{fig:Figure1}) is given 
	by \citep[e.g.,][]{bhattacharyya1967,gunn1975}:
	\begin{equation}
		\mathcal{F} [ d^{o}(\mathbf{r}^{\dagger}) ] =
		\alpha_{d} \, \mathcal{F} [ U(\mathbf{r}^{\dagger}) ] \: ,
		\quad \mathbf{r}^{\dagger} \in \mathcal{S}^{\dagger} \: ,
		\label{eq:observed-data-Fourier}
	\end{equation}
	where 
	\begin{equation}
		\mathcal{F} [ U(\mathbf{r}^{\dagger}) ] = 2 \pi \kappa 
		\int\limits_{z^{t}_{0}}^{\infty}
		\mathcal{F} [ \sigma(\mathbf{r}') ] \:
		\frac{e^{- k \, (z' - z^{\dagger}_{0})}}{k}
		\: dz' \: , 
		\quad \mathbf{r}^{\dagger} \in \mathcal{S}^{\dagger} \: ,
		\label{eq:potential-U-Fourier}
	\end{equation}
	is the Fourier transform of the Newtonian potential 
	$U$ (equation \ref{eq:potential-U-space}) on $\mathcal{S}^{\dagger}$, with $k = \sqrt{k_{x}^{2} + k_{y}^{2}}$ ,
	\begin{equation}
		\mathcal{F} [ \sigma(\mathbf{r}') ] = 
		\int\limits_{-\infty}^{\infty}\int\limits_{-\infty}^{\infty}
		\sigma(\mathbf{r}') \: e^{-\mathtt{i} \, (k_{x} \, x' + k_{y} \, y')}
		\: dx'dy'
		\label{eq:sigma-Fourier}
	\end{equation}
	is the Fourier transform of $\sigma(\mathbf{r}')$ 
	(the sources physical property) on the horizontal plane at the vertical
	coordinate $z'$ and $\mathtt{i}$ represents the imaginary unit.
	In equation \ref{eq:observed-data-Fourier}, $\alpha_{d}$ is a function of the 
	wavenumbers $k_{x}$ and $k_{y}$ and represents the differential operator 
	$D_{d}$ (equation \ref{eq:observed-data-space}) in the Fourier domain.
	
	\subsection{Differential operators}
	\label{sec:diff_operators}
	
	The observed potential field $d^{o}$ is represented in both the space and Fourier 
	domains using the differential operators $D_{d}$ and $\alpha_{d}$, as defined in equations 
	\ref{eq:observed-data-space} and \ref{eq:observed-data-Fourier}, respectively. 
	These operators act on the Newtonian potential $U$, described by equations 
	\ref{eq:potential-U-space} and \ref{eq:potential-U-Fourier}. To support the subsequent analysis, 
	it is essential to establish a relationship that enables the derivation of the Fourier-domain 
	operator $\alpha_{d}$ from its space-domain counterpart $D_{d}$ and vice-versa.  
	A possible relationship can be given by
	\begin{equation}
		D(\ell, \hat{\mathbf{u}}) = 
		\left( \hat{\mathbf{u}}^{\top} \nabla \right)^{\ell} \Leftrightarrow
		\alpha(\ell, \hat{\mathbf{u}}) = \left( \hat{\mathbf{u}}^{\top} \mathbf{k} \right)^{\ell} \: ,
		\quad \ell \in \{ 0, 1, 2 \} \: ,
		\label{eq:diff-operators-generic}
	\end{equation}
	where $\hat{\mathbf{u}}$ is a constant unit vector,
	$\nabla$ is the vector differential operator nabla \citep[e.g.,][ p. 179]{kaplan2002} and 
	$\mathbf{k} = ( \mathtt{i} \, k_{x} \, , \, \mathtt{i} \, k_{y} \, , \, k)^{\top}$
	\citep{pedersen1978}.
	The differential operator $D$ (equation \ref{eq:diff-operators-generic})
	has a symbolic form based on that used by \citet[][ p. 179]{kaplan2002}.
	
	By combining successive applications of $D(\ell, \hat{\mathbf{u}})$ (equation \ref{eq:diff-operators-generic})
	with properly chosen $\ell$ and $\hat{\mathbf{u}}$, we may define differential operators that 
	can be used to define not only the observed potential field $d^{o}$, but also any other potential 
	field commonly used in Geophysics.
	Some of them are shown in Table \ref{tbl:Table2}.
	Note that $D(0, \hat{\mathbf{u}}) = \alpha(0, \hat{\mathbf{u}}) = 1$ produces no differentiation.
	It can be verified, for example, that using $\ell = 1$ and a constant unit vector $\hat{\mathbf{u}}$
	in equation \ref{eq:diff-operators-generic} leads to the following directional derivative 
	of first order
	\begin{equation}
		D(1, \hat{\mathbf{u}}) \, U(\mathbf{r}) = \nabla U(\mathbf{r})^{\top} \hat{\mathbf{u}}
		\label{eq:potential-U-oblique-1st-deriv-space}
	\end{equation}
	where
	\begin{equation}
		\nabla U(\mathbf{r}) = \left( 
		\tfrac{\partial}{\partial x} U(\mathbf{r}),
		\tfrac{\partial}{\partial y} U(\mathbf{r}),
		\tfrac{\partial}{\partial z} U(\mathbf{r}) \right)^{\top}
		\label{eq:potential-U-grad-space}
	\end{equation}
	is the gradient of $U(\mathbf{r})$ (equation \ref{eq:potential-U-space}), with elements
	given by the partial derivatives of first order
	\begin{equation}
		\tfrac{\partial}{\partial u} U(\mathbf{r}) = \kappa \iiint\limits_{\mathcal{V}'}
		\sigma(\mathbf{r}') \: 
		\tfrac{\partial}{\partial u} \left( \frac{1}{\| \mathbf{r} - \mathbf{r}'\|} \right)
		\: d\mathcal{V}' \: , 
		\quad \mathbf{r} \notin \overline{\mathcal{V}'} \: , \quad u \in \{ x, y, z \} \: .
		\label{eq:potential-U-1st-deriv-space}
	\end{equation}
	The first-order directional derivative (equation \ref{eq:potential-U-oblique-1st-deriv-space})
	coincides with the partial derivatives of first order $\tfrac{\partial}{\partial u} U$ (equation \ref{eq:potential-U-1st-deriv-space})
	for the particular case in which the unit vector $\hat{\mathbf{u}}$ coincides with one of the canonical base vectors
	\begin{equation}
		\hat{\mathbf{x}} = (1,0,0)^{\top} , \quad 
		\hat{\mathbf{y}} = (0,1,0)^{\top} , \quad
		\hat{\mathbf{z}} = (0,0,1)^{\top} \: .
		\label{eq:canonical-base-vector}
	\end{equation}
	For example, $D(1, \hat{\mathbf{z}}) \, U(\mathbf{r}) = \tfrac{\partial}{\partial z} U(\mathbf{r})$.
	
	It can also be verified from equation \ref{eq:diff-operators-generic} that 
	\begin{equation}
		D(1, \hat{\mathbf{u}}) \, D(1, \hat{\mathbf{v}}) \, U(\mathbf{r}) = 
		\hat{\mathbf{u}}^{\top} \mathbf{H}_{U}(\mathbf{r}) \, \hat{\mathbf{v}} \: ,
		\label{eq:potential-U-oblique-2nd-deriv-space}
	\end{equation}
	represents a second-order directional derivative along constant unit vectors $\hat{\mathbf{u}}$ and
	$\hat{\mathbf{v}}$, where 
	\begin{equation}
		\mathbf{H}_{U}(\mathbf{r}) = 
		\begin{bmatrix}
			\tfrac{\partial^{2}}{\partial x^{2}} U(\mathbf{r}) & 
			\tfrac{\partial^{2}}{\partial x \, \partial y} U(\mathbf{r}) & 
			\tfrac{\partial^{2}}{\partial x \, \partial z} U(\mathbf{r}) \\
			\tfrac{\partial^{2}}{\partial x \, \partial y} U(\mathbf{r}) & 
			\tfrac{\partial^{2}}{\partial y^{2}} U(\mathbf{r}) & 
			\tfrac{\partial^{2}}{\partial y \, \partial z} U(\mathbf{r}) \\
			\tfrac{\partial^{2}}{\partial x \, \partial z} U(\mathbf{r}) & 
			\tfrac{\partial^{2}}{\partial y \, \partial z} U(\mathbf{r}) & 
			\tfrac{\partial^{2}}{\partial z^{2}} U(\mathbf{r})
		\end{bmatrix}
		\label{eq:potential-U-Hessian}
	\end{equation}
	is the Hessian matrix \citep[e.g.,][ p. 158]{kaplan2002} of $U$ (equation \ref{eq:potential-U-space}),
	with elements given by the partial derivatives of second order
	\begin{equation}
		\tfrac{\partial^{2}}{\partial u \, \partial v} U(\mathbf{r}) = \kappa \iiint\limits_{\mathcal{V}'}
		\sigma(\mathbf{r}') \: 
		\tfrac{\partial^{2}}{\partial u \, \partial v} \left( \frac{1}{\| \mathbf{r} - \mathbf{r}'\|} \right)
		\: d\mathcal{V}' \: , \quad \mathbf{r} \notin \overline{\mathcal{V}'} \: , \quad u, v \in \{ x, y, z \} \: .
		\label{eq:potential-U-2nd-deriv-space}
	\end{equation} 
	It can be easily verified that 
	the second-order directional derivative $D(1, \hat{\mathbf{u}}) \, D(1, \hat{\mathbf{v}}) \, U(\mathbf{r})$ 
	(equation \ref{eq:potential-U-oblique-2nd-deriv-space})
	coincides with the partial derivatives of second order $\tfrac{\partial^{2}}{\partial u \, \partial v} U(\mathbf{r})$ 
	(equation \ref{eq:potential-U-2nd-deriv-space}) for the particular cases in which the unit vectors 
	$\hat{\mathbf{u}}$ and $\hat{\mathbf{v}}$ are given by the canonical base vectors (equation \ref{eq:canonical-base-vector}).
	For example, $D(1, \hat{\mathbf{x}}) \, D(1, \hat{\mathbf{y}}) \, U(\mathbf{r}) = \tfrac{\partial^{2}}{\partial x \, \partial y} U(\mathbf{r})$.
	
	\subsection{Some relevant potential fields for geophysical applications}
	\label{subsec:relevant-potential-fields}
	
	Here we discuss some relevant potential fields for geophysical applications. 
	They are all listed in Table \ref{tbl:Table2}. 
	In gravity exploration methods, the vertical derivative $\tfrac{\partial}{\partial z} U$
	(equation \ref{eq:potential-U-1st-deriv-space}) does not define the gravity anomaly, but the \textit{gravity disturbance} data 
	\cite[e.g.,][ p. 93]{hofmann-wellenhof-moritz2005}.
	The debate around the differences between gravity anomaly and gravity disturbance is beyond the scope of the 
	present work and has been carried about from a more geodetic than a geophysical point of view 
	\citep[e.g.,][]{li2001, fairhead2003, hackney-featherstone2003, hinze2005, vajda-etal2006, vajda-etal2007, vajda-etal2008}. 
	Maybe, that is why it is still unfamiliar to a considerable group of exploration geophysicists.
	Another important quantity for gravity exploration methods is the Hessian matrix $\mathbf{H}_{U}$
	(equation \ref{eq:potential-U-Hessian}), which defines the
	(gravitational) \textit{gradient tensor} \citep[e.g.,][]{pedersen_rasmussen1990}.
	
	In magnetic exploration methods, the second-order directional derivative 
	$\hat{\mathbf{u}}^{\top} \mathbf{H}_{U} \, \hat{\mathbf{v}}$ (equation \ref{eq:potential-U-oblique-2nd-deriv-space})
	is commonly used to approximate the \textit{total-field anomaly} produced by 3D sources \cite[e.g., ][ p. 179]{blakely1996},
	where $\hat{\mathbf{u}}$ defines the constant direction of the main geomagnetic
	field at the study area and $\hat{\mathbf{v}}$ the constant total-magnetization direction of the sources.
	In this case, we use the term \textit{approximated total-field anomaly} to denote this 
	potential field throughout the article.
	Note that $\hat{\mathbf{u}}^{\top} \mathbf{H}_{U} \, \hat{\mathbf{v}}$ 
	(equation \ref{eq:potential-U-oblique-2nd-deriv-space}) can also be used to define (i) the  
	\textit{Cartesian components of the magnetic induction field} produced due to the 3D
	sources by setting the unit vector $\hat{\mathbf{u}}$ equal to the canonical base vectors $\hat{\mathbf{x}}$,
	$\hat{\mathbf{y}}$ and $\hat{\mathbf{z}}$ (equation \ref{eq:canonical-base-vector}) or (ii)
	the \textit{RTP anomaly} by setting $\hat{\mathbf{u}} = \hat{\mathbf{v}} = \pm\hat{\mathbf{z}}$.
	This special anomaly represents the vertical component of 
	the magnetic induction field that would be produced by the 3D sources if they had a 
	purely induced and vertical magnetization. Differently from generic total-field anomalies 
	\cite[e.g.,][ p. 178]{blakely1996}, which have a ``skewed'' pattern 
	oscillating from negative to positive values and do not have their maxima placed right 
	above the magnetic sources, the RTP anomaly is easier to interpret 
	because it is predominantly positive right above the sources and
	decays to zero towards their boundaries. \citet{baranov1957} was the first
	to propose a method, called \textit{reduction to the pole} (RTP), to transform generic
	total-field anomalies into this special anomaly.
	Finally, another common field in magnetic exploration methods is the amplitude of 
	the magnetic induction field generated by the true sources, defined by 
	$\| \nabla \left[\nabla U(\mathbf{r})^{\top} \hat{\mathbf{v}}\right] \|_{2}$. 
	This field is commonly known in applied geophysics as the 
	\textit{amplitude of the anomalous magnetic field} \citep{li_etal2010} or 
	\textit{amplitude of the magnetic anomaly vector} \citep{hidalgo_gato_etal2021, melo_etal2021} 
	and it is commonly used in magnetic interpretation as an alternative to the RTP anomaly 
	because it does not require information about the total-magnetization direction of the 
	true sources.
	
	
	\subsection{Classical Dirichlet's and Neumann's problems on a plane}
	\label{subsec:dirichlet_neumann}
	
	Due to the properties imposed on $\sigma(\mathbf{r}')$ (Table \ref{tbl:Table1}), the potential 
	$U$ (equation \ref{eq:potential-U-space}) has continuous partial derivatives of all orders, 
	satisfies the Laplace's equation at all points $\mathbf{r} \notin \overline{\mathcal{V}'}$ 
	located outside the sources, as well as is regular at infinity \citep[][ p. 122--124, 217]{kellogg1967}.
	Because of these properties, we have that
	\begin{equation}
		U(\mathbf{r}) = 
		\int\limits_{-\infty}^{\infty} \int\limits_{-\infty}^{\infty}
		D(\ell, \hat{\mathbf{z}}) U(\mathbf{r}'') \:\: 
		D(1 - \ell, \hat{\mathbf{z}}) \Psi(\mathbf{r}, \mathbf{r}'')
		\: d\mathcal{S}'' \: , \quad  z < z''_{0} < z^{t}_{0} \: ,
		\quad \ell \in \{0, 1\} \: ,
		\label{eq:U-Dirichlet-Neumann}
	\end{equation}
	where $\mathbf{r} = (x, y, z)^{\top} \notin \overline{\mathcal{V}'}$; 
	$\mathbf{r}'' = (x'', y'', z''_{0})^{\top} \in \mathcal{S}''$ (Figure \ref{fig:Figure1}); 
	$D(\ell, \hat{\mathbf{z}})$ is the space-domain differential operator 
	defined by equation \ref{eq:diff-operators-generic} in terms of the canonical base vector 
	$\hat{\mathbf{z}}$ (equation \ref{eq:canonical-base-vector}) and 
	\begin{equation}
		\Psi(\mathbf{r}, \mathbf{r}'') = 
		\frac{1}{2\pi} \: 
		\frac{1}{\| \mathbf{r} - \mathbf{r}'' \|} \: , 
		\quad \mathbf{r} \ne \mathbf{r}'' \: ,
		\label{eq:Psi}
	\end{equation}
	represents the scaled inverse distance function between $\mathbf{r}$ and $\mathbf{r}''$. 
	Equation \ref{eq:U-Dirichlet-Neumann} is commonly deduced from 
	the \textit{Green's third identity} \citep[][ p. 219]{kellogg1967}, according to
	the general outline presented by \cite{roy1962}, for example.
	Figure \ref{fig:Figure1} illustrates all geometric elements that are relevant in equations  
	\ref{eq:U-Dirichlet-Neumann} and \ref{eq:Psi}.
	
	For $\ell = 0$, equation \ref{eq:U-Dirichlet-Neumann} states that $U$ can be defined
	in terms of its values $D(0, \hat{\mathbf{z}}) U(\mathbf{r}'') = U(\mathbf{r}'')$ on a horizontal surface $\mathcal{S}''$ 
	above the sources and represents the solution of the classical \textit{Dirichlet’s problem} or the 
	\textit{first boundary value problem of potential theory} \cite[][ p. 236]{kellogg1967} on a plane, where 
	$D(1, \hat{\mathbf{z}}) \Psi(\mathbf{r}, \mathbf{r}'') = \tfrac{\partial}{\partial z} \Psi(\mathbf{r}, \mathbf{r}'')$ 
	represents the Green's function.
	From the physical point of view, equation \ref{eq:U-Dirichlet-Neumann} with $\ell = 0$ defines a continuous 
	layer of monopoles (or vertically-magnetized dipoles) that is located on $\mathcal{S}''$, has 
	physical-property distribution $U(\mathbf{r}'')$ and exactly reproduces the 
	potential field $U(\mathbf{r})$ at any point $\mathbf{r}$ above $\mathcal{S}''$.
	
	On the other hand, equation \ref{eq:U-Dirichlet-Neumann} for $\ell = 1$ shows that $U$ can also
	be defined in terms of its vertical derivative 
	$D(1, \hat{\mathbf{z}}) U(\mathbf{r}'') = \tfrac{\partial}{\partial z} U(\mathbf{r}'')$ 
	on a horizontal surface $\mathcal{S}''$ above the sources and represents the solution of the classical \textit{Neumann’s problem} or the 
	\textit{second boundary value problem of potential theory} \cite[][ p. 246]{kellogg1967} on a plane, 
	with the Green's function given by $D(0, \hat{\mathbf{u}}) \Psi(\mathbf{r}, \mathbf{r}'') = \Psi(\mathbf{r}, \mathbf{r}'')$.
	Equation \ref{eq:U-Dirichlet-Neumann} with $\ell = 1$ can be alternatively interpreted as a planar layer of
	vertically-magnetized dipoles (or monopoles) that is located on $\mathcal{S}''$, has 
	physical-property distribution $\tfrac{\partial}{\partial z} U(\mathbf{r}'')$ and exactly reproduces 
	$U(\mathbf{r})$ at any point $\mathbf{r}$ above $\mathcal{S}''$.
	
	By applying the Fourier-convolution theorem \cite[][ p. 268]{blakely1996} to equation \ref{eq:U-Dirichlet-Neumann} 
	on a horizontal surface $\mathcal{S}^{\dagger}$ above the sources, we find that
	\begin{equation}
		\mathcal{F} [ U(\mathbf{r}^{\dagger}) ] = 
		k \: \mathcal{F} [ U(\mathbf{r}'') ] \:
		\mathcal{F} [ \Psi(\mathbf{r}^{\dagger}, \mathbf{r}'') ] \: ,
		\quad z^{\dagger}_{0} < z''_{0} < z^{t}_{0} \: , \quad k \ne 0 \: ,
		\label{eq:U-Dirichlet-Neumann-Fourier}
	\end{equation}
	where $\mathbf{r}^{\dagger} = (x^{\dagger}, y^{\dagger}, z^{\dagger}_{0})^{\top} \in \mathcal{S}^{\dagger}$; 
	$\mathbf{r}'' = (x'', y'', z''_{0})^{\top} \in \mathcal{S}''$ (Figure \ref{fig:Figure1}) and
	\begin{equation}
		\mathcal{F} [ \Psi(\mathbf{r}^{\dagger}, \mathbf{r}'') ] 
		= \frac{e^{- k \, (z''_{0} - z^{\dagger}_{0})}}{k}
		\: , \quad z^{\dagger}_{0} < z''_{0} \: , \quad k \ne 0 \: ,
		\label{eq:Psi-Fourier}
	\end{equation}
	is defined according to \citet[][ p. 275]{blakely1996}.
	Figure \ref{fig:Figure1} illustrates all geometric elements that are relevant in equations  
	\ref{eq:U-Dirichlet-Neumann-Fourier} and \ref{eq:Psi-Fourier}.
	We get from equation \ref{eq:diff-operators-generic} that the term $k$ in equation \ref{eq:U-Dirichlet-Neumann-Fourier} represents,
	in the Fourier domain, the combined effect of the space-domain differential operators $D(\ell, \hat{\mathbf{z}})$ and 
	$D(1-\ell, \hat{\mathbf{z}})$.
	
	\subsection{Potential-field transformations}
	\label{subsec:potential_field_transformations}
	
	Several geophysical applications use the observed potential field 
	$d^{o}$ (equation \ref{eq:observed-data-space}) to estimate another potential field 
	defined in the space domain by
	\begin{equation}
		f(\mathbf{r}) = D_{f} \, U(\mathbf{r}) \: , 
		\quad \mathbf{r} \notin \overline{\mathcal{V}'} \: ,
		\label{eq:transformed-field-space}
	\end{equation}
	or, in the Fourier domain, 
	\begin{equation}
		\mathcal{F} [ f(\mathbf{r}^{\dagger}) ] =
		\alpha_{f} \, \mathcal{F} [ U(\mathbf{r}^{\dagger}) ] \: ,
		\quad \mathbf{r}^{\dagger} \in \mathcal{S}^{\dagger} \: ,
		\label{eq:transformed-field-Fourier}
	\end{equation}
	where $\mathcal{S}^{\dagger}$ is a horizontal surface defined above the sources (Figure \ref{fig:Figure1});
	$D_{f}$ is a space-domain differential operator acting on $U$ (equation \ref{eq:potential-U-space}) and
	$\alpha_{f}$ its corresponding representation in the Fourier domain.
	The differential operators $D_{f}$ and $\alpha_{f}$ associated with the potential field 
	$f$ (equations \ref{eq:transformed-field-space} and \ref{eq:transformed-field-Fourier}) are commonly
	defined in terms of $D$ and $\alpha$ (equation \ref{eq:diff-operators-generic}). 
	Estimating $f$ (equations \ref{eq:transformed-field-space} and \ref{eq:transformed-field-Fourier}) from 
	the observed potential field $d^{o}$ (equation \ref{eq:observed-data-space}) is an operation 
	commonly known in geophysics as \textit{potential-field transformation}.
	
	A common approach to compute potential-field transformations \cite[e.g.,][]{gunn1975} 
	consists in using the definition of observed data in the Fourier domain (equation \ref{eq:observed-data-Fourier}) 
	and manipulating the Fourier domain solution of the Dirichlet's and Neumann's problems 
	(equation \ref{eq:U-Dirichlet-Neumann-Fourier}) to obtain
	\begin{equation}
		\mathcal{F} [ f(\mathbf{r}^{\dagger}) ] =
		\left( \tfrac{\alpha_{f}}{\alpha_{d}} \right) \, \mathcal{F} \left[ d^{o}(\mathbf{r}^{\dagger}) \right] \: , 
		\quad \mathbf{r}^{\dagger} \in \mathcal{S}^{\dagger} \: .
		\label{eq:field-transformation-Fourier}
	\end{equation}
	Equation \ref{eq:field-transformation-Fourier} defines a single-step potential-field transformation 
	in terms of the observed potential field $d^{o}(\mathbf{r}^{\dagger})$ on a planar surface $\mathcal{S}^{\dagger}$ 
	located above the sources, at $z^{\dagger}_{0}$ (Figure \ref{fig:Figure1}). 
	The potential field $f$ in the space domain is obtained by simply computing the inverse Fourier 
	transform of \ref{eq:field-transformation-Fourier}.
	For simplicity, both fields $d^o$ and $f$ are defined at the same points in equation \ref{eq:field-transformation-Fourier}, 
	although it's worth noting that, in certain geophysical contexts, $f$ may be estimated at different locations from 
	where potential-field data were measured.
	
	The \textit{equivalent-layer technique} offers an alternative approach that is generally defined in the 
	space domain to compute potential-field transformations. What follows aims at discussing the well-established 
	empirical approaches used in the equivalent-layer technique, as well as discussing the theoretical background 
	supporting them.
	
	% Planar equivalent layers
	\section{Planar equivalent layers}
	\label{eq:planar_eqlayer}
	
	An equivalent layer is loosely defined as a physical property distribution that satisfactorily 
	reproduces potential-field data and enables potential-field transformations. 
	In the literature, equivalent layers are typically assumed to lie in source-free regions where 
	Laplace’s equation holds. Consequently, it is presumed that the equivalent layer does not intersect 
	the true sources. However, in practical applications, the positions of the true sources in the subsurface 
	are often unknown, making this condition difficult to guarantee.
	Empirical evidence suggests that equivalent-layer methods can still perform effectively 
	when the layer intersects the sources \citep[e.g.,][]{gonzalez_etal2022}.
	
	In this section, we formally define \textit{discrete equivalent layers} and \textit{continuous equivalent layers} 
	on a horizontal surface $\mathcal{S}''$ that may either lie above or intersect the sources (Figure \ref{fig:Figure1}). 
	Notably, any discrete equivalent layer can be constructed by discretizing a continuous equivalent layer.
	Our goal is to leverage the properties of continuous layers to gain deeper insights into their discrete counterparts, 
	which will be explored in detail in subsequent sections.
	
	\begin{definition}[\textbf{Discrete equivalent layer}]
		Let us assume to have an $N \times 1$ observed data vector $\mathbf{d}^{o}$ whose elements 
		are the potential-field data $d^{o}_{i}$ at observation points $\mathbf{r}_{i}$ outside the sources, 
		$i \in \{1, \dots, N\}$, above a horizontal surface $\mathcal{S}''$ with constant vertical 
		coordinate $z''_{0}$. Besides, consider an $N \times M$ 
		matrix $\mathbf{A}$ with element $a_{ij}$ defined by a harmonic function that depends only on $\mathbf{r}_{i}$ 
		and $\mathbf{r}''_{j} \in \mathcal{S}''$, $j \in \{1, \dots, M\}$, and
		vanishes as the distance $\| \mathbf{r}_{i} - \mathbf{r}''_{j} \|_{2}$ increases. 
		Then, for a chosen accuracy $\epsilon > 0$, we may call the $M \times 1$ vector $\mathbf{p}_{\epsilon}$ with $j$-th element $p_{j}$ 
		a \underline{discrete equivalent layer} if its predicted potential field 
		\begin{equation}
			\mathbf{d}_{\epsilon} = \mathbf{A} \, \mathbf{p}_{\epsilon}
			\label{eq:predicted-data-discrete-layer}
		\end{equation}
		satisfies 
		\begin{equation}
			\| \mathbf{d}^{o} - \mathbf{d}_{\epsilon} \|_{2}^{2} \le \epsilon \: .
			\label{eq:epsilon-bound-data-discrete}
		\end{equation}
		\label{def:discrete-eqlayer}
	\end{definition}
	
	In Definition \ref{def:discrete-eqlayer}, $\mathbf{d}_{\epsilon}$ is the predicted data vector whose 
	$i$-th element represents the potential field due to $\mathbf{p}_{\epsilon}$ at the observation point $\mathbf{r}_{i}$. 
	Besides, the elements of $\mathbf{A}$ are usually defined by 
	$a_{ij} \equiv D_{a} \Psi(\mathbf{r}_{i}, \mathbf{r}''_{j})$, which represents a 
	partial derivative of $\Psi$ (equation \ref{eq:Psi}), with $D_{a}$ given by the space-domain differential operator 
	$D$ (equation \ref{eq:diff-operators-generic}). 
	It is important to bear in mind that the horizontal surface $\mathcal{S}''$ in Definition \ref{def:discrete-eqlayer} 
	may cross the sources, as illustrated in Figure \ref{fig:Figure1}.
	The main assumption in Definition \ref{def:discrete-eqlayer} is that $d^{o}_{i}$ can be approximated 
	within a given accuracy $\epsilon$ by the potential field produced by a discrete and finite set of $M$ 
	virtual sources such as monopoles (or point masses) or dipoles with 
	constant and unknown physical property $p_{j}$ (e.g., mass, magnetic-moment intensity). 
	These virtual sources are usually called \textit{equivalent sources}.
	We emphasize that Definition \ref{def:discrete-eqlayer} presumes a planar equivalent layer formed 
	by punctual equivalent sources that satisfies the data fit criterion given by equation \ref{eq:epsilon-bound-data-discrete}, 
	but these premises are certainly open to adjustments if necessary.
	
	After estimating a particular $\mathbf{p}_{\epsilon}$ (Definition \ref{def:discrete-eqlayer}), 
	a potential field transformation such as upward continuation, 
	interpolation or reduction to the pole, for example, can be computed as follows:
	\begin{equation}
		\mathbf{f}_{\epsilon} = \mathbf{T} \, \mathbf{p}_{\epsilon} \: ,
		\label{eq:transformed-field-discrete-layer} 
	\end{equation}
	where $\mathbf{f}_{\epsilon}$ is an $L \times 1$ vector whose $k$-th element approximates 
	the transformed potential field $f(\mathbf{r}^{\ddagger}_{k})$ (equation \ref{eq:transformed-field-space}) 
	at a point $\mathbf{r}^{\ddagger}_{k} \in \mathcal{S}^{\ddagger}$ above $\mathcal{S}''$ (Figure \ref{fig:Figure1})
	and $\mathbf{T}$ is an $L \times M$ matrix whose element $t_{kj} \equiv D_{t} \Psi(\mathbf{r}^{\ddagger}_{k}, \mathbf{r}''_{j})$,
	$k \in \{ 1, \dots, L \}$, $j \in \{ 1, \dots, M \}$, 
	is commonly defined by a partial derivative of $\Psi$ (equation \ref{eq:Psi}), with $D_{t}$ given by the 
	space-domain differential operator $D$ (equation \ref{eq:diff-operators-generic}).
	The element $t_{kj}$ represents the transformed potential field due to the $j$-th equivalent source. 
	We have defined the transformed potential field at points 
	$\mathbf{r}^{\ddagger}_{k} \in \mathcal{S}^{\ddagger}$ above $\mathcal{S}''$ (Figure \ref{fig:Figure1}), however they 
	can be defined at any place any place above $\mathcal{S}''$, including the observation points. 
	
	From the physical point of view, the product of the $j$-th element $p_{j}$ of $\mathbf{p}_{\epsilon}$ 
	and the harmonic function $a_{ij}$ representing the element $ij$ of matrix $\mathbf{A}$ 
	(equation \ref{eq:predicted-data-discrete-layer}) defines the potential field produced at the $i$-th observation
	point by the $j$-th equivalent source with scalar physical property $p_{j}$.
	For example, if $\mathbf{d}^{o}$ contains gravity disturbance data
	\cite[e.g.,][ p. 93]{hofmann-wellenhof-moritz2005}, then
	$p_{j} \, a_{ij}$ commonly defines the vertical component of the gravitational attraction
	produced at $\mathbf{r}_{i}$ by a monopole located at
	$\mathbf{r}''_{j} \in \mathcal{S}''$ (Figure \ref{fig:Figure1}), with mass $p_{j}$. 
	This is the pioneer approach presented by \cite{dampney1969} for 
	processing gravity data.
	For the case in which $\mathbf{d}^{o}$ contains total-field anomaly data 
	\cite[e.g.,][ p. 178]{blakely1996} instead, the common approach consists in
	following \cite{emilia1973} and defining the term $p_{j} \, a_{ij}$ such
	that it represents a magnetic induction field component 
	produced at the $i$-th observation point by a dipole located at 
	$\mathbf{r}''_{j} \in \mathcal{S}''$ (Figure \ref{fig:Figure1}), with magnetic moment intensity $p_{j}$.
	
	Note that, in the examples presented above, the physical relationship 
	between the observed potential field data and their true sources is the
	same as that between the data and the equivalent sources.
	However, many authors 
	\cite[e.g.,][]{cordell1992, guspi_novara2009, li_etal2014, soler_uieda2021} 
	have proposed equivalent-layer methods by assuming without any proof that the harmonic functions $a_{ij}$ are 
	independent of the nature of the observed potential-field data and can be 
	conveniently chosen according to the characteristics of the problem.
	We mathematically confirm this independence some sections later and also show how 
	the particular form of predefined harmonic functions $a_{ij}$ affect the subsequent 
	potential-field transformations.
	
	\begin{definition}[\textbf{Continuous equivalent layer}]
		Let us assume to have an observed potential field $d^{o}(\mathbf{r})$ defined above 
		a horizontal surface $\mathcal{S}''$ with constant vertical coordinate $z''_{0}$. 
		Besides, consider (i) a horizontal surface $\mathcal{S}^{\dagger}$ that is located above 
		$\mathcal{S}''$, at a constant vertical coordinate $z^{\dagger}_{0}$, 
		and does not cross the sources and (ii) a harmonic kernel 
		$a(\mathbf{r}, \mathbf{r}'')$ that depends only on $\mathbf{r}$ and $\mathbf{r}'' \in \mathcal{S}''$, as well as 
		vanishes as the distance $\| \mathbf{r} - \mathbf{r}'' \|_{2}$ increases. 
		Then, for a chosen accuracy $\varepsilon > 0$, we may call $p_{\varepsilon}(\mathbf{r}'')$ 
		a \underline{continuous equivalent layer} if its predicted potential field 
		\begin{equation}
			d_{\varepsilon}(\mathbf{r})
			= \int\limits_{-\infty}^{\infty}\int\limits_{-\infty}^{\infty}
			p_{\varepsilon}(\mathbf{r}'') \, a(\mathbf{r}, \mathbf{r}'') \:\: d\mathcal{S}''
			\: , \quad z < z''_{0} \: ,
			\label{eq:predicted-data-continuous-layer}
		\end{equation}
		satisfies 
		\begin{equation}
			\int\limits_{-\infty}^{\infty} \int\limits_{-\infty}^{\infty}
			\vert d^{o}(\mathbf{r}^{\dagger}) - 
			d_{\varepsilon}(\mathbf{r}^{\dagger}) \vert^{2}
			\:\: d\mathcal{S}^{\dagger} \le \varepsilon \: ,
			\quad z^{\dagger}_{0} < z^{t}_{0} \: , \quad z^{\dagger}_{0} < z''_{0} \: ,
			\label{eq:epsilon-bound-data}
		\end{equation}
		where $\mathbf{r} = (x, y, z)^{\top}$, 
		$\mathbf{r}^{\dagger} = (x^{\dagger}, y^{\dagger}, z^{\dagger}_{0})^{\top} \in \mathcal{S}^{\dagger}$ 
		and $\mathbf{r}'' = (x'', y'', z''_{0})^{\top} \in \mathcal{S}''$.
		\label{def:continuous-eqlayer}
	\end{definition}
	
	In Definition \ref{def:continuous-eqlayer}, $p_{\varepsilon}(\mathbf{r}'')$ represents a continuous 
	physical-property distribution that extends to infinity along the horizontal directions. 
	Moreover, the kernel function $a(\mathbf{r}, \mathbf{r}'')$ is usually defined by $D_{a} \Psi(\mathbf{r}, \mathbf{r}'')$ 
	as a generic partial derivative of $\Psi$ (equation \ref{eq:Psi}),
	where $D_{a}$ is set in terms of the space-domain differential operator 
	$D$ (equation \ref{eq:diff-operators-generic}).
	It is important to stress here that the horizontal surface $\mathcal{S}''$ in Definition \ref{def:continuous-eqlayer} 
	may cross the sources, according to Figure \ref{fig:Figure1}.
	
	Similarly to the discrete case (equation \ref{eq:transformed-field-discrete-layer}), a particular continuous equivalent layer 
	$p_{\varepsilon}(\mathbf{r}'')$ leads to a transformed potential field given by:
	\begin{equation}
		f_{\varepsilon}(\mathbf{r}^{\ddagger})
		= \int\limits_{-\infty}^{\infty}\int\limits_{-\infty}^{\infty}
		p_{\varepsilon}(\mathbf{r}'') \, t(\mathbf{r}^{\ddagger}, \mathbf{r}'') \:\: d\mathcal{S}'' \: , 
		\quad \mathbf{r}^{\ddagger} \in \mathcal{S}^{\ddagger} \: ,
		\quad \mathbf{r}'' \in \mathcal{S}'' \: ,
		\label{eq:transformed-field-continuous-layer}
	\end{equation}
	where $\mathbf{r}^{\ddagger}$ are points defined on a planar surface $\mathcal{S}^{\ddagger}$ located above $\mathcal{S}''$ (Figure \ref{fig:Figure1}); 
	$f_{\varepsilon}(\mathbf{r}^{\ddagger})$ is a transformed field that approximates 
	$f(\mathbf{r}^{\ddagger})$ (equations \ref{eq:transformed-field-space}); and
	$t(\mathbf{r}^{\ddagger}, \mathbf{r}'')$ is the transformation kernel,
	which is usually defined by $D_{t} \Psi(\mathbf{r}^{\ddagger}, \mathbf{r}'')$ as a generic partial 
	derivative of the potential function $\Psi$
	(equation \ref{eq:Psi}), with $D_{t}$ given in terms of $D$ (equation \ref{eq:diff-operators-generic}).
	In equation \ref{eq:transformed-field-continuous-layer}, we define $\mathbf{r}^{\ddagger}$ on 
	a planar surface $\mathcal{S}^{\ddagger}$ located above $\mathcal{S}''$ (Figure \ref{fig:Figure1}), but this is not a 
	restriction. They can coincide with the observation points or be arbitrarily defined above $\mathcal{S}''$, for example.
	
	% Planar equivalent-layer theorem
	\section{Existence of continuous equivalent layers on a plane}
	\label{sec:planar_eqlayer_theorem}
	
	\subsection{Modified Dirichlet's and Neumann's problems on a plane}
	
	The classical solutions to Dirichlet’s and Neumann’s problems (equations \ref{eq:U-Dirichlet-Neumann} 
	and \ref{eq:U-Dirichlet-Neumann-Fourier}) on a plane are specifically formulated for cases where 
	the horizontal surface $\mathcal{S}''$ is positioned above the sources. In this section, we extend 
	these classical solutions by removing this restriction and allowing $\mathcal{S}''$ to intersect 
	the sources (Figure \ref{fig:Figure1}).
	
	The modified Dirichlet’s and Neumann’s problems introduced here, 
	which accommodate horizontal surfaces intersecting the sources, represent a pivotal step in our work. 
	They establish the mathematical foundation necessary for rigorously proving the existence of planar equivalent layers.
	
	We start by considering an auxiliary potential $U_{\varepsilon}$ that approximates $U$ (equation \ref{eq:potential-U-space})
	at points $\mathbf{r} \notin \overline{\mathcal{V}'}$ above the sources, is square integrable on the surface 
	$\mathcal{S}''$ and has a space domain representation given by 
	\begin{equation}
		U_{\varepsilon}(\mathbf{r}) = 
		\int\limits_{-\infty}^{\infty} \int\limits_{-\infty}^{\infty}
		D(\ell, \hat{\mathbf{z}}) U_{\varepsilon}(\mathbf{r}'') \:\:
		D(1-\ell, \hat{\mathbf{z}}) \Psi(\mathbf{r}, \mathbf{r}'') 
		\:\: d\mathcal{S}'' \: , 
		\quad z < z^{t}_{0} \: , \quad z < z''_{0} \: ,
		\quad \ell \in \{ 0, 1 \} \: ,
		\label{eq:potential-U-epsilon-space}
	\end{equation}
	where $\mathbf{r} = (x, y, z)^{\top} \notin \overline{\mathcal{V}'}$; 
	$\mathbf{r}'' = (x'', y'', z''_{0})^{\top} \in \mathcal{S}''$ (Figure \ref{fig:Figure1}) and 
	$D$ is the space-domain differential operator 
	defined by equation \ref{eq:diff-operators-generic} in terms of the canonical base vector 
	$\hat{\mathbf{z}}$ (equation \ref{eq:canonical-base-vector}).
	We also consider that the Fourier-convolution theorem 
	can be applied to the auxiliary potential $U_{\varepsilon}$
	(equation \ref{eq:potential-U-epsilon-space}) on the horizontal surface $\mathcal{S}^{\dagger}$ above
	the sources (Figure \ref{fig:Figure1}), so that
	\begin{equation}
		\mathcal{F} [ U_{\varepsilon}(\mathbf{r}^{\dagger}) ]
		= k \, \mathcal{F} [ U_{\varepsilon}(\mathbf{r}'') ] \:
		\mathcal{F} [ \Psi(\mathbf{r}^{\dagger}, \mathbf{r}'') ] \: ,
		\quad \mathbf{r}^{\dagger} \in \mathcal{S}^{\dagger} \: .
		\label{eq:potential-U-epsilon-Fourier-dagger}
	\end{equation}
	Using the definition of $\mathcal{F} [ \Psi(\mathbf{r}^{\dagger}, \mathbf{r}'') ]$ (equation \ref{eq:Psi-Fourier})
	in equation \ref{eq:potential-U-epsilon-Fourier-dagger}, we conclude that $U_{\varepsilon}(\mathbf{r}'')$ 
	must satisfy
	\begin{equation}
		\mathcal{F} [ U_{\varepsilon}(\mathbf{r}'') ] = 
		\mathcal{F} [ U_{\varepsilon}(\mathbf{r}^{\dagger}) ] \, e^{\, k \, (z''_{0} - z^{\dagger}_{0})} \: .
		\label{eq:potential-U-epsilon-Fourier-layer}
	\end{equation}
	Now we present two possible approaches to define an auxiliary potential $U_{\varepsilon}$ that satisfies the 
	conditions above and lead to modified Dirichlet’s and Neumann’s problems that can be applied to 
	the case in which $\mathcal{S}''$ cross the sources (Figure \ref{fig:Figure1}). 
	
	\subsubsection{Approach 1: Truncating wavenumber}
	
	Let $U_{\varepsilon}$ be given by
	\begin{equation}
		\mathcal{F} [ U_{\varepsilon}(\mathbf{r}^{\dagger}) ] = 
		\begin{cases}
			\mathcal{F} [ U(\mathbf{r}^{\dagger}) ] &
			, \quad 0 < k \le k_{\varepsilon} \\
			0 & , \quad k > k_{\varepsilon} 
		\end{cases} \: ,
		\label{eq:potential-U-epsilon-Fourier-dagger-1}
	\end{equation}
	where $k_{\varepsilon}$ is a fixed and possibly very large truncating wavenumber and 
	$\mathcal{F} \left[ U(\mathbf{r}^{\dagger}) \right]$ is defined on the horizontal surface $\mathcal{S}^{\dagger}$ 
	located above the sources (Figure \ref{fig:Figure1}). 
	As we can see, this definition of $U_{\varepsilon}$ 
	(i) imposes that its Fourier transform coincides with 
	that of $U$ (equation \ref{eq:potential-U-Fourier}) on the horizontal surface $\mathcal{S}^{\dagger}$ 
	up to the truncating wavenumber $k_{\varepsilon}$; 
	(ii) represents a smooth version of $U$ on $\mathcal{S}^{\dagger}$; and
	(iii) leads to a square integrable function on $\mathcal{S}^{\dagger}$ because 
	$U$ is square integrable on $\mathcal{S}^{\dagger}$ (Appendix A).
	Then we obtain from equation \ref{eq:potential-U-epsilon-Fourier-layer} that 
	\begin{equation}
		\mathcal{F} \left[ U_{\varepsilon}(\mathbf{r}'') \right] = 
		\begin{cases}
			\mathcal{F} \left[ U(\mathbf{r}^{\dagger}) \right] 
			e^{\, k \, (z''_{0} - z^{\dagger}_{0})} & , 
			\quad 0 < k \le k_{\varepsilon} \\
			0 & , \quad k > k_{\varepsilon} 
		\end{cases} \: ,
		\label{eq:potential-U-epsilon-Fourier-layer-1}
	\end{equation}
	which approximates the potential field obtained by 
	computing the downward continuation of $U$ (equation \ref{eq:potential-U-space}) 
	to points $\mathbf{r}'' \in \mathcal{S}''$ that may be located inside the sources (Figure \ref{fig:Figure1}). 
	It is well-known, however, that downward continuation can only be computed at 
	source-free points, where the Laplace's equation applies. This means that the
	continuation plane is limited by the top $z^{t}_{0}$ of the sources, otherwise the continued field 
	presents an oscillatory character \citep[e.g.,][]{roy1967, fedi_florio2002} that may
	prevent $\mathcal{F} \left[ U_{\varepsilon}(\mathbf{r}'') \right]$ to be square integrable.
	The truncating wavenumber $k_{\varepsilon}$ used to define $U_{\varepsilon}$ according to equation 
	\ref{eq:potential-U-epsilon-Fourier-layer-1} overcomes this problem, as suggested (but not shown) 
	by \citet{bhattacharyya1967},
	and guarantees that the Fourier-convolution theorem can be applied to equation \ref{eq:potential-U-epsilon-space} 
	for obtaining equation \ref{eq:potential-U-epsilon-Fourier-dagger}.
	
	At last, it can be easily verified that, if we let the truncating wavenumber $k_{\varepsilon}$ 
	goes to infinity in equation \ref{eq:potential-U-epsilon-Fourier-layer-1} for the case in which 
	$\mathcal{S}''$ is located above the sources, at $z''_{0} < z^{t}_{0}$ (Figure \ref{fig:Figure1}), 
	we get 
	\begin{equation}
		\lim_{k_{\varepsilon}\to\infty} \mathcal{F} \left[ U_{\varepsilon}(\mathbf{r}'') \right] = 
		\mathcal{F} \left[ U(\mathbf{r}'') \right] \: , \quad z''_{0} < z^{t}_{0} \: ,
		\label{eq:lim-kepsilon-potential-U-epsilon-Fourier-layer}
	\end{equation}
	\begin{equation}
		\lim_{k_{\varepsilon}\to\infty} \mathcal{F} \left[ U_{\varepsilon}(\mathbf{r}^{\dagger}) \right] = 
		\mathcal{F} \left[ U(\mathbf{r}^{\dagger}) \right] \: , \quad z^{\dagger}_{0} < z''_{0} < z^{t}_{0} \: ,
		\label{eq:lim-kepsilon-potential-U-epsilon-Fourier-dagger}
	\end{equation}
	pointwise and, consequently, that
	\begin{equation}
		\lim_{k_{\varepsilon}\to\infty} U_{\varepsilon}(\mathbf{r}) = 
		U(\mathbf{r}) \: , \quad z < z''_{0} < z^{t}_{0} \: ,
		\label{eq:lim-kepsilon-potential-U-epsilon-space}
	\end{equation}
	which is in perfect agreement with the classical solutions of the Dirichlet's and Neumann's
	problems (equations \ref{eq:U-Dirichlet-Neumann} and \ref{eq:U-Dirichlet-Neumann-Fourier}).
	
	\subsubsection{Approach 2: Tikhonov regularization}
	
	We may also define $U_{\varepsilon}$ as the following smooth approximation of $U$ on the horizontal surface 
	$\mathcal{S}^{\dagger}$ located above the sources (Figure \ref{fig:Figure1}):
	\begin{equation}
		\mathcal{F} [ U_{\varepsilon}(\mathbf{r}^{\dagger}) ] = 
		\tfrac{1}{1 + \mu_{\varepsilon} \, k^{2} \, e^{\, k \, (z''_{0} - z^{\dagger}_{0})}} \; 
		\mathcal{F} [ U(\mathbf{r}^{\dagger}) ] \: ,
		\label{eq:potential-U-epsilon-Fourier-dagger-2}
	\end{equation}
	where $\mu_{\varepsilon}$ is a fixed positive regularization parameter imposing smoothness on $U$ 
	according to \citet{pasteka_etal2012}. 
	Again, this proposed definition of $U_{\varepsilon}$ 
	(i) represents a smooth version of $U$ on $\mathcal{S}^{\dagger}$ and
	(ii) leads to a square integrable function on $\mathcal{S}^{\dagger}$.
	Next, we obtain from equation \ref{eq:potential-U-epsilon-Fourier-layer} that 
	\begin{equation}
		\mathcal{F} \left[ U_{\varepsilon}(\mathbf{r}'') \right] = 
		\tfrac{e^{\, k \, (z''_{0} - z^{\dagger}_{0})}}{1 + \mu_{\varepsilon} \, k^{2} \, e^{\, k \, (z''_{0} - z^{\dagger}_{0})}} \; 
		\mathcal{F} [ U(\mathbf{r}^{\dagger}) ] \: .
		\label{eq:potential-U-epsilon-Fourier-layer-2}
	\end{equation}
	The alternative definition of $U_{\varepsilon}$ given by equation \ref{eq:potential-U-epsilon-Fourier-dagger-2} 
	also leads to an approximation of the potential field obtained by 
	computing the downward continuation of $U$ (equation \ref{eq:potential-U-space}) 
	to points $\mathbf{r}'' \in \mathcal{S}''$ that may be located inside the sources (Figure \ref{fig:Figure1}). 
	In this case, however, it is the regularization parameter $\mu_{\varepsilon}$ that forces  
	$U_{\varepsilon}$ to be square integrable on $\mathcal{S}''$ and guarantees that the 
	Fourier-convolution theorem can be applied to equation \ref{eq:potential-U-epsilon-space} 
	for obtaining equation \ref{eq:potential-U-epsilon-Fourier-dagger}.
	
	Finally, if we let the $\mu_{\varepsilon}$ 
	goes to zero in equation \ref{eq:potential-U-epsilon-Fourier-layer-2} for the case in which 
	$\mathcal{S}''$ is located above the sources, at $z''_{0} < z^{t}_{0}$ (Figure \ref{fig:Figure1}), 
	we obtain once more that 
	\begin{equation}
		\lim_{\mu_{\varepsilon}\to 0} \mathcal{F} \left[ U_{\varepsilon}(\mathbf{r}'') \right] = 
		\mathcal{F} \left[ U(\mathbf{r}'') \right] \: , \quad z''_{0} < z^{t}_{0} \: ,
		\label{eq:lim-muepsilon-potential-U-epsilon-Fourier-layer}
	\end{equation}
	\begin{equation}
		\lim_{\mu_{\varepsilon}\to 0} \mathcal{F} \left[ U_{\varepsilon}(\mathbf{r}^{\dagger}) \right] = 
		\mathcal{F} \left[ U(\mathbf{r}^{\dagger}) \right] \: , \quad z^{\dagger}_{0} < z''_{0} < z^{t}_{0} \: ,
		\label{eq:lim-muepsilon-potential-U-epsilon-Fourier-dagger}
	\end{equation}
	pointwise and, consequently, that
	\begin{equation}
		\lim_{\mu_{\varepsilon}\to 0} U_{\varepsilon}(\mathbf{r}) = 
		U(\mathbf{r}) \: , \quad z < z''_{0} < z^{t}_{0} \: ,
		\label{eq:lim-muepsilon-potential-U-epsilon-space}
	\end{equation}
	which is also in perfect agreement with the classical solutions of the Dirichlet's and Neumann's
	problems (equations \ref{eq:U-Dirichlet-Neumann} and \ref{eq:U-Dirichlet-Neumann-Fourier}).
	
	\subsection{Formal existence demonstration of continuous layers}
	
	In this subsection, we formally demonstrate the existence of planar equivalent layers 
	$p_{\varepsilon}(\mathbf{r}'')$ (Definition \ref{def:continuous-eqlayer}) 
	through the formulation and proof of a theorem. This approach provides, for the first time, a rigorous mathematical 
	framework for a concept traditionally based on empirical assumptions.
	As the cornerstone of this study, the theorem moves beyond empirical assumptions to establish 
	a mathematically sound basis for the well-known applications of the equivalent-layer technique, 
	even for the case in which the layer crosses the sources.
	
	\begin{theorem}[\textbf{Planar equivalent-layer theorem}]
		\label{theorem:planar-eqlayer-theorem}
		Let a horizontal surface $\mathcal{S}''$ and a kernel function $a(\mathbf{r}, \mathbf{r}'')$ 
		be set according to Definition \ref{def:continuous-eqlayer}, for a given observed potential field $d^{o}$. 
		Then, for every $\varepsilon > 0$, there exists a set of continuous equivalent layers $p_{\varepsilon}$ on $\mathcal{S}''$. 
		Furthermore, for the case in which $\mathcal{S}''$ is above the sources, the continuous equivalent layers 
		reproduce $d^{o}$ exactly.
	\end{theorem}
	\begin{proof}
		We start by defining $p_{\varepsilon}$ (Definition \ref{def:continuous-eqlayer}) 
		in the Fourier domain via
		\begin{equation}
			\mathcal{F} [ p_{\varepsilon}(\mathbf{r}'') ] =
			\tfrac{k \, \alpha_{d}}{\alpha_{a}} \,
			\mathcal{F} [ U_{\varepsilon}(\mathbf{r}'') ] \: ,
			\quad \mathbf{r}'' \in \mathcal{S}'' \: ,
			\label{eq:p-Fourier}
		\end{equation}
		where $\alpha_{a}$ is a differential operator 
		defined in the Fourier domain in terms of $\alpha$ (equation \ref{eq:diff-operators-generic}) and
		$U_{\varepsilon}(\mathbf{r}'')$ is defined as in \ref{eq:potential-U-epsilon-Fourier-layer-1},
		for a given truncating wavenumber $k_{\varepsilon}$, or as in \ref{eq:potential-U-epsilon-Fourier-layer-2},
		for a given regularization parameter $\mu_{\varepsilon}$. 
		
		Note that there is an equivalent layer $p_{\varepsilon}$ for each particular $\alpha_{a}$.
		For any finite choice of $0 < k_{\varepsilon} < \infty$ or $0 < \mu_{\varepsilon} < \infty$
		(equations \ref{eq:potential-U-epsilon-Fourier-layer-1} and \ref{eq:potential-U-epsilon-Fourier-layer-2}), 
		both $\mathcal{F} [ p_{\varepsilon}(\mathbf{r}'') ]$ 
		and $p_{\varepsilon}(\mathbf{r}'')$ are square-integrable and, consequently, it is possible to
		apply the Fourier-convolution theorem to \ref{eq:predicted-data-continuous-layer} and obtain
		\begin{equation}
			\mathcal{F} [ d_{\varepsilon}(\mathbf{r}^{\dagger}) ] =
			\mathcal{F} [ p_{\varepsilon}(\mathbf{r}'') ] \, 
			\mathcal{F} [ a(\mathbf{r}^{\dagger}, \mathbf{r}'') ] \: ,
			\quad \mathbf{r}^{\dagger} \in \mathcal{S}^{\dagger} \: ,
			\quad \mathbf{r}'' \in \mathcal{S}'' \: ,
			\label{eq:predicted-data-continuous-layer-Fourier}
		\end{equation}
		where 
		\begin{equation}
			\mathcal{F} [ a(\mathbf{r}^{\dagger}, \mathbf{r}'') ] = \alpha_{a} \, \mathcal{F} [ \Psi(\mathbf{r}^{\dagger}, \mathbf{r}'') ] \: ,
			\quad \mathbf{r}^{\dagger} \in \mathcal{S}^{\dagger} \: ,
			\quad \mathbf{r}'' \in \mathcal{S}'' \: ,
			\label{eq:kernel-a-Fourier}
		\end{equation}
		with $a(\mathbf{r}, \mathbf{r}'') = D_{a} \Psi(\mathbf{r}, \mathbf{r}'')$;
		$\mathbf{r}^{\dagger} = (x^{\dagger}, y^{\dagger}, z^{\dagger}_{0})^{\top}$; 
		$\mathbf{r}'' = (x'', y'', z''_{0})^{\top}$ and $z^{\dagger}_{0} < z''_{0}$ (Table \ref{tbl:Table1}).
		By substituting equations \ref{eq:observed-data-Fourier}, \ref{eq:Psi-Fourier} and \ref{eq:kernel-a-Fourier} 
		into \ref{eq:predicted-data-continuous-layer-Fourier} and comparing the result with equations 
		\ref{eq:potential-U-epsilon-Fourier-layer-1} and \ref{eq:potential-U-epsilon-Fourier-layer-2}, 
		we find that $d_{\varepsilon}$ satisfies
		\begin{equation}
			\mathcal{F} [ d_{\varepsilon}(\mathbf{r}^{\dagger}) ] = 
			\alpha_{d} \,
			\mathcal{F} [ U_{\varepsilon}(\mathbf{r}^{\dagger}) ] \: ,
			\quad \mathbf{r}^{\dagger} \in \mathcal{S}^{\dagger} \: .
			\label{eq:predicted-data-continuous-layer-epsilon-Fourier}
		\end{equation}
		
		From Plancherel's theorem \citep[][ p. 119]{bracewell1999} and use of polar coordinates 
		$k_{x} = k \, \cos\theta$, $k_{y} = k \, \sin\theta$ and $k = \sqrt{k_{x}^{2} + k_{y}^{2}}$,
		we get
		\begin{equation}
			\begin{split}
				\int\limits_{-\infty}^{\infty} \int\limits_{-\infty}^{\infty}
				\vert 
				d^{o}(\mathbf{r}^{\dagger}) - d_{\varepsilon}(\mathbf{r}^{\dagger})
				\vert^{2} 
				\: d\mathcal{S}^{\dagger} = 
				\int\limits_{-\infty}^{\infty} \int\limits_{-\infty}^{\infty}
				\vert \mathcal{F} [ d^{o}(\mathbf{r}^{\dagger}) ] - \mathcal{F} [
				d_{\varepsilon}(\mathbf{r}^{\dagger}) ] \vert^{2} 
				\: d k_{x} \, d k_{y} \\
				= \int\limits_{0}^{2\pi} \int\limits_{0}^{\infty}
				\vert \mathcal{F} [ d^{o}(\mathbf{r}^{\dagger}) ] - \mathcal{F} [
				d_{\varepsilon}(\mathbf{r}^{\dagger}) ] \vert^{2} \: k \: d k \, d \theta \: .
			\end{split}
			\label{eq:Plancherel}
		\end{equation}
		If we consider that the auxiliary potential $U_{\varepsilon}$ is defined by equation \ref{eq:potential-U-epsilon-Fourier-dagger-1},
		we obtain that 
		\begin{equation}
			\mathcal{F} [ d_{\varepsilon}(\mathbf{r}^{\dagger}) ] = 
			\begin{cases}
				\mathcal{F} [ d^{o}(\mathbf{r}^{\dagger}) ] &
				, \quad 0 < k \le k_{\varepsilon} \\
				0 & , \quad k > k_{\varepsilon} 
			\end{cases} \: .
			\label{eq:predicted-data-continuous-layer-epsilon-Fourier-1}
		\end{equation}
		If we now consider that $U_{\varepsilon}$ is given by equation \ref{eq:potential-U-epsilon-Fourier-dagger-2},
		we get
		\begin{equation}
			\mathcal{F} [ d_{\varepsilon}(\mathbf{r}^{\dagger}) ] = 
			\tfrac{1}{1 + \mu_{\varepsilon} \, k^{2} \, e^{\, k \, (z''_{0} - z^{\dagger}_{0})}} \; 
			\mathcal{F} [ d^{o}(\mathbf{r}^{\dagger}) ] \: .
			\label{eq:predicted-data-continuous-layer-epsilon-Fourier-2}
		\end{equation}
		It is important to note that using a finite but sufficiently large $0 < k_{\varepsilon} < \infty$ in equation 
		\ref{eq:predicted-data-continuous-layer-epsilon-Fourier-1} or a finite $0 < \mu_{\varepsilon} < \infty$ in 
		equation \ref{eq:predicted-data-continuous-layer-epsilon-Fourier-2} leads to  
		$\int_{0}^{2\pi} \int_{0}^{\infty} \vert \mathcal{F} [ d^{o}(\mathbf{r}^{\dagger}) ] - \mathcal{F} [d_{\varepsilon}(\mathbf{r}^{\dagger}) ] \vert^{2} \: k \, dkd\theta \le \varepsilon$ in equation \ref{eq:Plancherel}. 
		Because of that, we conclude that $p_{\varepsilon}$ (equation \ref{eq:p-Fourier})
		is a continuous equivalent layer according to Definition \ref{def:continuous-eqlayer}. 
		Besides, if the layer plane is above the sources (i.e., if $z''_{0} < z^{t}_{0}$), then both 
		$\mathcal{F} [ p_{\varepsilon}(\mathbf{r}'') ]$ and $p_{\varepsilon}(\mathbf{r}'')$ are square-integrable
		even for the choice $k_{\varepsilon} = \infty$ or $\mu_{\varepsilon} = 0$. Then, in that case, 
		$\int_{0}^{2\pi} \int_{0}^{\infty} 
		\vert \mathcal{F} [ d^{o}(\mathbf{r}^{\dagger}) ] - \mathcal{F} [d_{\varepsilon}(\mathbf{r}^{\dagger}) ] \vert^{2} \: k \, dkd\theta 
		= 0$ and $p_{\varepsilon}$ yields a predicted field $d_{\varepsilon}$ that exactly reproduces the observed field $d^{o}$ 
		on $\mathcal{S}^{\dagger}$.
	\end{proof}
	
	\subsubsection{Remark 1}
	
	The planar equivalent-layer theorem guarantees the existence of a planar and continuous equivalent layer, even if
	it crosses the sources, i.e., if $z''_{0} > z^{t}_{0}$ (Figure \ref{fig:Figure1}).
	But it does not provide any guidance on the particular choice of the 
	truncating wavenumber $k_{\varepsilon}$ (equation \ref{eq:potential-U-epsilon-Fourier-dagger-1}) or the 
	regularization parameter $\mu_{\varepsilon}$ (equation \ref{eq:potential-U-epsilon-Fourier-dagger-2}) 
	for the construction of $p_{\varepsilon}$, with a given accuracy level 
	$\varepsilon$ (Definition \ref{def:continuous-eqlayer}).
	
	The absence of a rule to construct $k_{\varepsilon}$ (equation \ref{eq:potential-U-epsilon-Fourier-dagger-1}) or 
	$\mu_{\varepsilon}$ (equation \ref{eq:potential-U-epsilon-Fourier-dagger-2}) in the 
	planar equivalent-layer theorem (Theorem \ref{theorem:planar-eqlayer-theorem}) may be a little disconcerting 
	at first sight, but this is very common in mathematical analysis. 
	Our theorem is a nonconstructive existence statement \citep[e.g., ][ p. 200]{tao_2016}, such as the well-known
	\textit{intermediate value theorem} from calculus, for example.
	
	\subsubsection{Remark 2}
	
	The truncating wavenumber $k_{\varepsilon}$ and the regularization parameter $\mu_{\varepsilon}$ 
	play a central role in the planar equivalent-layer theorem (Theorem \ref{theorem:planar-eqlayer-theorem}). 
	Both parameters impose smoothness on the planar equivalent layer. 
	From a theoretical point of view, it might be interesting to study the influence of
	discretizing continuous equivalent layers and its relation to $k_{\varepsilon}$ and $\mu_{\varepsilon}$. 
	In practical situations, the interpreter usually estimate a ``sufficiently smooth'' discrete equivalent layer 
	that ``satisfactorily'' fit the observed data within a given error level. 
	It seems that this discretization and the imposed degree of smoothness are important factors rendering 
	the equivalent-layer technique even when the layer plane crosses the sources.
	We implicitly assume that the estimated discrete layer $\mathbf{p}_{\epsilon}$ (Definition \ref{def:discrete-eqlayer}) 
	approximates the corresponding continuous layer $p_{\varepsilon}$ (Definition \ref{def:continuous-eqlayer}) 
	as the discretization becomes increasingly finer. 
	This assumption is confirmed by the numerical tests we present later.
	
	\subsubsection{Remark 3}
	
	Any continuous equivalent layer $p_{\varepsilon}$ (Definition \ref{def:continuous-eqlayer}) leads to a predicted 
	potential field that satisfies
	the error bound $\varepsilon$ (equation \ref{eq:epsilon-bound-data}) not only on the surface $\mathcal{S}^{\dagger}$, but also 
	on any $\mathcal{S}^{\ddagger}$ with $z^{\ddagger}_{0} < z^{\dagger}_{0}$ (Figure \ref{fig:Figure1}). 
	Namely, it follows from equation \ref{eq:potential-U-Fourier} that
	\begin{equation}
		\int\limits_{-\infty}^{\infty} \int\limits_{-\infty}^{\infty}
		\vert d^{o}(\mathbf{r}^{\ddagger}) - d_{\varepsilon}(\mathbf{r}^{\ddagger})
		\vert^{2} \: d\mathcal{S}^{\ddagger} \le
		\int\limits_{-\infty}^{\infty} \int\limits_{-\infty}^{\infty}
		\vert d^{o}(\mathbf{r}^{\dagger}) - d_{\varepsilon}(\mathbf{r}^{\dagger})
		\vert^{2} \: d\mathcal{S}^{\dagger} \: ,
		\quad z^{\ddagger}_{0} < z^{\dagger}_{0} \: .
		\label{eq:error-bound-upper-planes}
	\end{equation}
	Thus, if $p_{\varepsilon}$ produces a $d_{\varepsilon}$ satisfying 
	$\int_{-\infty}^{\infty} \int_{-\infty}^{\infty} \vert d^{o}(\mathbf{r}^{\dagger}) - d_{\varepsilon}(\mathbf{r}^{\dagger}) 
	\vert^{2} \: d\mathcal{S}^{\dagger} \le \varepsilon$, then it automatically holds that 
	$\int_{-\infty}^{\infty} \int_{-\infty}^{\infty} \vert d^{o}(\mathbf{r}^{\ddagger}) - d_{\varepsilon}(\mathbf{r}^{\ddagger}) 
	\vert^{2} \: d\mathcal{S}^{\ddagger} \le \varepsilon$.
	
	% Potential-field transformations
	\section{Unified approach for potential-field transformations}
	
	Given a particular equivalent layer $p_{\varepsilon}$ (Definition \ref{def:continuous-eqlayer}) 
	defined as a square-integrable function by equation \ref{eq:p-Fourier}, in terms of specific $\alpha_{d}$ and $\alpha_{a}$, 
	it is theoretically possible to compute any potential field transformation (equations \ref{eq:transformed-field-space} and 
	\ref{eq:transformed-field-Fourier}).
	It means that, contrarily to the common sense, we do not need to use a planar layer of dipoles to compute the RTP anomaly from 
	the approximated total-field anomaly, for example. In fact, even a layer of monopoles can be used to perform this potential-field transformation. 
	This can be verified by first noting that applying the Fourier-convolution theorem to \ref{eq:transformed-field-continuous-layer} 
	on a planar surface $\mathcal{S}^{\ddagger}$ (Figure \ref{fig:Figure1}) gives
	\begin{equation}
		\mathcal{F} [ f_{\varepsilon}(\mathbf{r}^{\ddagger}) ] =
		\mathcal{F} \left[ p_{\varepsilon}(\mathbf{r}'') \right]
		\mathcal{F} [ t(\mathbf{r}^{\ddagger}, \mathbf{r}'') ] \: , 
		\quad \mathbf{r}^{\ddagger} \in \mathcal{S}^{\ddagger} \: ,
		\quad \mathbf{r}'' \in \mathcal{S}'' \: ,
		\label{eq:transformed-field-continuous-layer-Fourier}
	\end{equation}
	where
	\begin{equation}
		\mathcal{F} [ t(\mathbf{r}^{\ddagger}, \mathbf{r}'') ] = \alpha_{t} \, \mathcal{F} [ \Psi(\mathbf{r}^{\ddagger}, \mathbf{r}'') ]
		\label{eq:kernel-t-Fourier}
	\end{equation}
	is defined in terms of a differential operator $\alpha_{t}$ acting on $\Psi$ (equation \ref{eq:Psi}) in the Fourier domain.
	Then we set 
	\begin{equation}
		\alpha_{t} = \left( \tfrac{\alpha_{a}}{\alpha_{d}} \right) \alpha_{f} \: ,
		\label{eq:alpha-t-generic}
	\end{equation}
	where $\alpha_{f}$ is the differential operator associated with the desired potential field $f$ (equation \ref{eq:transformed-field-Fourier}).
	In this case, we get from equations \ref{eq:Psi-Fourier}, \ref{eq:p-Fourier}  and \ref{eq:kernel-t-Fourier} 
	that the transformed potential field $f_{\varepsilon}$ (equation \ref{eq:transformed-field-continuous-layer-Fourier}) 
	is given by:
	\begin{equation}
		\mathcal{F} [ f_{\varepsilon}(\mathbf{r}^{\ddagger}) ] =
		\alpha_{f} \, k \, \mathcal{F} [ U_{\varepsilon}(\mathbf{r}'') ] \: 
		\mathcal{F} [ \Psi(\mathbf{r}^{\ddagger}, \mathbf{r}'') ] \: , 
		\quad \mathbf{r}^{\ddagger} \in \mathcal{S}^{\ddagger} \: ,
		\quad \mathbf{r}'' \in \mathcal{S}'' \: .
		\label{eq:transformed-field-continuous-layer-Fourier2}
	\end{equation}
	Finally, we use \ref{eq:Psi-Fourier}, \ref{eq:transformed-field-Fourier} and \ref{eq:potential-U-epsilon-Fourier-layer} to show that
	\begin{equation}
		\mathcal{F} [ f_{\varepsilon}(\mathbf{r}^{\ddagger}) ] = 
		\begin{cases}
			\mathcal{F} [ f(\mathbf{r}^{\ddagger}) ] &
			, \quad 0 < k \le k_{\varepsilon} \\
			0 & , \quad k > k_{\varepsilon} 
		\end{cases} \: ,
		\label{eq:transformed-field-continuous-layer-1-Fourier}
	\end{equation}
	if the auxiliary potential $U_{\varepsilon}$ is defined by equation \ref{eq:potential-U-epsilon-Fourier-dagger-1},
	or
	\begin{equation}
		\mathcal{F} [ f_{\varepsilon}(\mathbf{r}^{\ddagger}) ] = 
		\tfrac{1}{1 + \mu_{\varepsilon} \, k^{2} \, e^{\, k \, (z''_{0} - z^{\dagger}_{0})}} \; 
		\mathcal{F} [ f(\mathbf{r}^{\ddagger}) ] \: ,
		\label{eq:transformed-field-continuous-layer-2-Fourier}
	\end{equation}
	if the auxiliary potential $U_{\varepsilon}$ is defined by equation \ref{eq:potential-U-epsilon-Fourier-dagger-2}.
	In both cases, we can see that the transformed field $f_{\varepsilon}$ yielded by the equivalent layer 
	is a smooth approximation of that produced by the true sources on the planar surface $\mathcal{S}^{\ddagger}$
	(Figure \ref{fig:Figure1}). Note that, to obtain a correct potential-field transformation leading to an $f_{\varepsilon}$ 
	that approximates $f$, we need to correctly set the operator $\alpha_{t}$ (equation \ref{eq:alpha-t-generic}), 
	regardless of the equivalent layer type defined by $\alpha_{a}$.
	
	
	In the context of the equivalent-layer technique, a potential-field transformation is generally defined
	in the space domain, by following two steps. The first consists in obtaining a particular layer $p_{\varepsilon}$ 
	(Definition \ref{def:continuous-eqlayer}), which is subsequently used to compute a transformed field $f_{\varepsilon}$ via 
	\ref{eq:transformed-field-continuous-layer}.
	It can be verified, however, that this two-step approach is equivalent to the single-step approach based on the 
	Fourier transform of the observed potential field (equation \ref{eq:field-transformation-Fourier}). 
	We find by using \ref{eq:observed-data-Fourier} and \ref{eq:transformed-field-continuous-layer-Fourier2} that
	\begin{equation}
		\mathcal{F} [ f_{\varepsilon}(\mathbf{r}^{\ddagger}) ] =
		\left( \tfrac{\alpha_{f}}{\alpha_{d}} \right) \, \mathcal{F} [ d_{\varepsilon}(\mathbf{r}^{\ddagger}) ]
		\: .
		\label{eq:transformed-field-continuous-layer-single-step}
	\end{equation}
	Then, if we consider that the predicted field $d_{\varepsilon}$ (equation \ref{eq:predicted-data-continuous-layer-epsilon-Fourier})
	is defined in terms of the auxiliary potential $U_{\varepsilon}$ given by equation \ref{eq:potential-U-epsilon-Fourier-dagger-1},
	we get
	\begin{equation}
		\mathcal{F} [ f_{\varepsilon}(\mathbf{r}^{\ddagger}) ] =
		\left( \tfrac{\alpha_{f}}{\alpha_{d}} \right) \, 
		\begin{cases}
			\mathcal{F} \left[ d^{o}(\mathbf{r}^{\dagger}) \right] 
			e^{\, k \, (z^{\ddagger}_{0} - z^{\dagger}_{0})} & , \quad 0 < k \le k_{\varepsilon} \\
			0 & , \quad k > k_{\varepsilon} 
		\end{cases} \: .
		\label{eq:transformed-field-continuous-layer-1-single-step}
	\end{equation}
	Alternatively, if we consider that the predicted field $d_{\varepsilon}$ (equation \ref{eq:predicted-data-continuous-layer-epsilon-Fourier})
	is defined in terms of the auxiliary potential $U_{\varepsilon}$ given by equation \ref{eq:potential-U-epsilon-Fourier-dagger-2},
	we get
	\begin{equation}
		\mathcal{F} [ f_{\varepsilon}(\mathbf{r}^{\ddagger}) ] =
		\left( \tfrac{\alpha_{f}}{\alpha_{d}} \right) \, 
		\tfrac{1}{1 + \mu_{\varepsilon} \, k^{2} \, e^{\, k \, (z''_{0} - z^{\dagger}_{0})}} \;
		\mathcal{F} \left[ d^{o}(\mathbf{r}^{\dagger}) \right] \: .
		\label{eq:transformed-field-continuous-layer-2-single-step}
	\end{equation}
	Note that, for the particular case in which $z^{\ddagger}_{0}$ coincides with $z^{\dagger}_{0}$ and 
	the truncating wavenumber $k_{\varepsilon}$ goes to infinity, equation 
	\ref{eq:transformed-field-continuous-layer-1-single-step} reduces to \ref{eq:field-transformation-Fourier}.
	Similarly, if $z^{\ddagger}_{0}$ coincides with $z^{\dagger}_{0}$ and $\mu_{\varepsilon}$ goes to zero,
	equation \ref{eq:transformed-field-continuous-layer-2-single-step} also reduces to \ref{eq:field-transformation-Fourier}.
	For the first time, we provide a mathematical proof establishing the link between
	potential-field transformations performed using the equivalent-layer technique in 
	the space domain and those computed via Fourier transformation of the observed data.
	This connection sheds light on why potential-field transformations based on the
	equivalent layers can exhibit the same numerical instabilities commonly associated
	with Fourier-based transformations of potential-field data. Although this
	relationship is widely acknowledged, it is typically attributed to the notion that
	physical phenomena are independent of the domain -- whether spatial or Fourier --
	used for their computation.
	
	% Equivalent-layer classes
	\section{Equivalent layer classes}
	\label{sec:eqlayer_classes}
	
	According to the planar equivalent-layer theorem (Theorem \ref{theorem:planar-eqlayer-theorem}), 
	it is possible to define different planar equivalent layers $p_{\varepsilon}(\mathbf{r}'')$, 
	$\mathbf{r}'' \in \mathcal{S}''$ (Figure \ref{fig:Figure1}), 
	satisfying the conditions imposed by Definition \ref{def:continuous-eqlayer}.
	Specifically, different equivalent layers $p_{\varepsilon}(\mathbf{r}'')$ can be defined
	on the same horizontal surface $\mathcal{S}''$, for the same potential field data $d^{o}$, 
	by using different combinations of $\alpha_{d}$ and $\alpha_{a}$ in equation \ref{eq:p-Fourier}. 
	The horizontal surface $\mathcal{S}''$ may be located above the sources 
	or intersect them (Figure \ref{fig:Figure1}).
	Here we show different classes of equivalent layers, discuss their applicability and define how they 
	can be used to perform potential-field transformations.
	
	\subsection{Dirichlet's equivalent layers}
	
	\begin{definition}
		An equivalent layer $p_{\varepsilon}$ (Definition \ref{def:continuous-eqlayer}) is called Dirichlet's equivalent 
		layer if it has the following form in the Fourier domain
		\begin{equation}
			\mathcal{F} [ p_{\varepsilon}(\mathbf{r}'') ] =
			\alpha_{d} \, \mathcal{F} [ U_{\varepsilon}(\mathbf{r}'') ]  \: .
			\label{eq:p-Fourier-Dirichlet}
		\end{equation}
		\label{def:eql-Dirichlet}
	\end{definition}
	Note that Dirichlet's layers (Definition \ref{def:eql-Dirichlet}) are obtained from 
	equation \ref{eq:p-Fourier} by setting $\alpha_{a} \equiv \alpha(1, \hat{\mathbf{z}}) = k$ 
	(equation \ref{eq:diff-operators-generic}).
	In this case, we can easily verify from equations \ref{eq:observed-data-Fourier},
	\ref{eq:potential-U-epsilon-Fourier-layer-1} and \ref{eq:potential-U-epsilon-Fourier-layer-2} that they 
	approximate the downward continuation of the observed potential field $d^{o}$ to the layer plane $\mathcal{S}''$.
	It is important noting that such layers depend on $\alpha_{d}$, which 
	carries information about the observed potential field (equation \ref{eq:observed-data-Fourier}).
	In this case, the corresponding predicted field (equation \ref{eq:predicted-data-continuous-layer}) assumes the form
	\begin{equation}
		d_{\varepsilon}(\mathbf{r})
		= \int\limits_{-\infty}^{\infty}\int\limits_{-\infty}^{\infty}
		p_{\varepsilon}(\mathbf{r}'') \: \tfrac{\partial}{\partial z} \Psi(\mathbf{r}, \mathbf{r}'') \:\: d\mathcal{S}''
		\: , \quad z < z''_{0} \: .
		\label{eq:predicted-data-Dirichlet-layer}
	\end{equation}
	From the physical point of view, this $d_{\varepsilon}$ is analogous
	to the vertical component of the gravitational attraction due to a layer of monopoles
	or the magnetic scalar potential due to a layer of vertically-magnetized dipoles.
	Moreover, equation \ref{eq:predicted-data-Dirichlet-layer} approximates the well-known \textit{upward continuation integral} 
	\citep[e.g.,][]{skeels1947, peters1949, henderson-zietz1949, dean1958, roy1962, bhattacharyya1967, henderson1970},
	which is traditionally derived from the solution of the Dirichlet’s problem on the
	layer plane (equation \ref{eq:U-Dirichlet-Neumann} with $\ell = 0$). 
	
	Dirichlet's equivalent layers are the basis of the pioneer method proposed by \cite{dampney1969} for processing gravity data.
	Such layers lead to potential-field transformations (equations 
	\ref{eq:transformed-field-continuous-layer} and \ref{eq:transformed-field-continuous-layer-Fourier}) defined by 
	$\alpha_{t} = \nicefrac{k \, \alpha_{f}}{\alpha_{d}}$ (equation \ref{eq:alpha-t-generic}).
	Dirichlet's equivalent layers are commonly used to perform upward/downward continuation of the observed potential 
	field, which is obtained by setting $\alpha_{f} = \alpha_{d}$, resulting in $\alpha_{t} = k$. 
	Other potential-field transformations can also be accomplished from such equivalent layers, but they must be 
	defined in terms of $\alpha_{f} \ne \alpha_{d}$.
	In this case, however, defining the resulting transformed field $f_{\varepsilon}$ explicitly in the space domain 
	(equation \ref{eq:transformed-field-continuous-layer}) 
	from its counterpart in the Fourier domain (equation \ref{eq:transformed-field-continuous-layer-Fourier}) may be very difficult.
	
	
	\subsection{Equivalent layers having minimum phase changes}
	
	\begin{definition}
		An equivalent layer $p_{\varepsilon}$ (Definition \ref{def:continuous-eqlayer}) is said to have minimum 
		phase changes (with respect to the spatial distribution of the true sources) if its phase spectrum is given by
		\begin{equation}
			\Theta [ p_{\varepsilon}(\mathbf{r}'') ] = \Theta [ U_{\varepsilon}(\mathbf{r}^{\dagger}) ] \: ,
			\label{eq:p-phase-spectrum-U}
		\end{equation}
		where $\Theta [ U_{\varepsilon}(\mathbf{r}^{\dagger}) ]$ denotes the phase spectrum of $U_{\varepsilon}$ 
		(Appendix B, equations \ref{eq:potential-U-epsilon-1-phase-spectrum} and \ref{eq:potential-U-epsilon-2-phase-spectrum}).
		\label{def:eql-min-phase}
	\end{definition}
	\begin{proposition}
		\label{prop:minimum-phase-change-layers}
		An equivalent layer $p_{\varepsilon}$ having minimum phase changes (Definition \ref{def:eql-min-phase})
		depends only on the spatial distribution of the true sources and not on the particular observed 
		potential field $d^{o}$.
	\end{proposition}
	\begin{proof}
		It follows from equations \ref{eq:p-phase-spectrum-U} and definition of phase spectrum (Appendix B) that 
		$\mathcal{F} [ p_{\varepsilon}(\mathbf{r}'') ]$ must be the product of a real function and
		$\mathcal{F} [ U_{\varepsilon}(\mathbf{r}^{\dagger}) ]$.
		From equations \ref{eq:potential-U-epsilon-Fourier-layer} and \ref{eq:p-Fourier}, we can verify that 
		$\mathcal{F} [ p_{\varepsilon}(\mathbf{r}'') ]$ is in fact the product of 
		$\mathcal{F} [ U_{\varepsilon}(\mathbf{r}^{\dagger}) ]$ 
		and a term $\left( \nicefrac{\alpha_{d}}{\alpha_{a}} \right) k \, e^{\, k \, (z''_{0} - z^{\dagger}_{0})}$. 
		Since $\nicefrac{\alpha_{d}}{\alpha_{a}}$ is the only term that can be a complex function, we conclude that 
		the imposed condition can only be satisfied for the case in which $\alpha_{a} = \alpha_{r} \, \alpha_{d}$, 
		where $\alpha_{r}$ is a generic real function satisfying $\alpha_{r} \ne 0$ for $0 < k \le k_{\varepsilon}$.
		It can then be verified that, in this case, \ref{eq:p-Fourier} assumes the form
		\begin{equation}
			\mathcal{F} [ p_{\varepsilon}(\mathbf{r}'') ] =
			\tfrac{k}{\alpha_{r}} \, \mathcal{F} [ U_{\varepsilon}(\mathbf{r}^{\dagger}) ] 
			e^{\, k \, (z''_{0} - z^{\dagger}_{0})} \: ,
			\label{eq:p-Fourier-min-phase}
		\end{equation}
		which is the product of a real function and $\mathcal{F} [ U_{\varepsilon}(\mathbf{r}^{\dagger}) ]$.
		Since $\alpha_{r}$ does not depend on $\alpha_{d}$, $\mathcal{F} [ p_{\varepsilon}(\mathbf{r}'') ]$ does not 
		depend on $\alpha_{d}$ as well. This implies that the phase spectrum of $p_{\varepsilon}$ approximates that 
		of the potential $U$ (Appendix B), thereby corresponding to an approximation 
		of the phase spectrum associated with the spatial distribution of the (shallow parts of the) true sources.
	\end{proof}
	
	Equivalent layers having minimum phase changes (Definition \ref{def:eql-min-phase}) may present intricate 
	spatial domain representations, dictated by the particular form of the real function $\alpha_{r}$
	(equation \ref{eq:p-Fourier-min-phase}).
	Such layers give raise to potential-field transformations (equation \ref{eq:transformed-field-continuous-layer-Fourier})
	defined by $\alpha_{t} = \alpha_{r} \, \alpha_{f}$ (equation \ref{eq:alpha-t-generic}), which may also
	have non-trivial representations in the space domain.
	
	\subsection{Neumann's equivalent layers}
	
	\begin{definition}
		An equivalent layer $p_{\varepsilon}$ (Definition \ref{def:continuous-eqlayer}) is called Neumann's equivalent 
		layer if it has the following form in the Fourier domain
		\begin{equation}
			\mathcal{F} [ p_{\varepsilon}(\mathbf{r}'') ] =
			k \, \mathcal{F} [ U_{\varepsilon}(\mathbf{r}'') ] \: .
			\label{eq:p-Fourier-Neumann}
		\end{equation}
		\label{def:eql-Neumann}
	\end{definition}
	\begin{proposition}
		All Neumann's equivalent layers have minimum phase changes.
		\label{prop:eql-Neumann}
	\end{proposition}
	\begin{proof}
		Equation \ref{eq:p-Fourier-Neumann} is directly obtained from \ref{eq:p-Fourier}, by using 
		$\alpha_{a} = \alpha_{d}$. It can be easily verified that, in this case, 
		$\mathcal{F} [ p_{\varepsilon}(\mathbf{r}'') ]$ is a particular case of equation \ref{eq:p-Fourier-min-phase},
		with $\alpha_{r} = 1$. Consequently, Neumann's equivalent layers have minimum phase changes
		according to Definition \ref{def:eql-min-phase}.
	\end{proof}
	
	The predicted potential field $d_{\varepsilon}$ (equation \ref{eq:predicted-data-continuous-layer}) 
	due to Neumann's equivalent layers (Definition \ref{def:eql-Neumann}) assumes the form
	\begin{equation}
		d_{\varepsilon}(\mathbf{r})
		= \int\limits_{-\infty}^{\infty}\int\limits_{-\infty}^{\infty}
		p_{\varepsilon}(\mathbf{r}'') \: D_{d} \Psi(\mathbf{r}, \mathbf{r}'') \:\: d\mathcal{S}''
		\: , \quad z < z''_{0} \: ,
		\label{eq:predicted-data-Neumann-layer}
	\end{equation}
	where $D_{d}$ is the same space-domain differential operator associated with
	the observed potential field data $d^{o}$ (equation \ref{eq:observed-data-space}).
	As a consequence, Neumann's equivalent layers (Definition \ref{def:eql-Neumann}) preserve the physical
	relationship between the observed potential field and the true sources.
	It means that, if the observed potential field $d^{o}$ is the $x$-component of the 
	gravitational attraction due to the true sources, the predicted potential field $d_{\varepsilon}$  
	represents the $x$-component of the gravitational attraction due to the planar equivalent layer, for example.
	The same is valid for any other potential field.
	
	Any particular Neumann's layer $p_{\varepsilon}$ yields transformed potential 
	fields (equation \ref{eq:transformed-field-continuous-layer-Fourier}) 
	defined by $\alpha_{t} = \alpha_{f}$ (equation \ref{eq:alpha-t-generic}) in the Fourier domain. 
	Hence, the corresponding transformed field $f_{\varepsilon}$ in the space domain 
	(equation \ref{eq:transformed-field-continuous-layer}) assumes the simple form
	\begin{equation}
		f_{\varepsilon}(\mathbf{r})
		= \int\limits_{-\infty}^{\infty}\int\limits_{-\infty}^{\infty}
		p_{\varepsilon}(\mathbf{r}'') \: D_{f} \Psi(\mathbf{r}, \mathbf{r}'') \:\: d\mathcal{S}''
		\: , \quad z < z''_{0} \: ,
		\label{eq:transformed-field-continuous-layer-Neumann}
	\end{equation}
	where $D_{f}$ is the space-domain representation of the differential operator $\alpha_{f}$ 
	associated with the desired transformed field $f$ (equation \ref{eq:transformed-field-space}).
	
	It follows immediately from equation \ref{eq:p-Fourier-Neumann} that Neumann's equivalent layers approximate 
	the vertical derivative $\tfrac{\partial}{\partial z} U(\mathbf{r}'')$ (equation \ref{eq:potential-U-1st-deriv-space}) 
	on the layer plane $\mathcal{S}''$.
	For the cases in which the observed potential field $d^{o}$ represents gravity data, 
	$p_{\varepsilon}$ approximates the gravity disturbance (or vertical component of the gravitational
	attraction) due to the true sources on the layer plane $\mathcal{S}''$.
	If $d^{o}$ represents magnetic data produced by sources with constant total magnetization direction, 
	then $p_{\varepsilon}$ approximates the pseudogravity 
	anomaly \citep[][ p. 345]{blakely1996} due to the true sources on $\mathcal{S}''$.
	In this case, $p_{\varepsilon}$ is strictly positive because it is defined in terms of a strictly positive 
	function $\sigma(\mathbf{r}')$ (equation \ref{eq:potential-U-space}) representing the 
	total-magnetization intensity at all points $\mathbf{r}' = (x', y', z')^{\top} \in \mathcal{V}'$ 
	within the sources.
	This result is in perfect agreement with that deduced by \citet{reis_etal2020} for
	the particular case in which the observed potential field $d^{o}$ is
	the approximated total-field anomaly (Table \ref{tbl:Table2}).
	
	It is widely known that integrating the potential field $\tfrac{\partial}{\partial z} U(\mathbf{r}'')$ 
	(equation \ref{eq:potential-U-1st-deriv-space}) on a horizontal plane above the causative sources
	gives their total mass (or total dipole moment), according to the
	Gauss' theorem \citep[e.g.,][]{hammer1945, lafehr1965, bhattacharyya1967}. 
	Furthermore, \citet{medeiros-silva1995a} show that the horizontal first-order moments, second-order moment $xy$, 
	difference between the second-order moments $xx$ and $yy$, as well as the horizontal coordinates of the sources center of mass
	(or center of dipole moment) of the true sources can also be estimated from the $\tfrac{\partial}{\partial z} U(\mathbf{r}'')$
	that they produce on a horizontal plane above them.
	Thus, all these parameters of the true sources can be directly estimated from 
	Neumann's equivalent layers $p_{\varepsilon}$ (Definition \ref{def:eql-Neumann}) because
	they approximate the field $\tfrac{\partial}{\partial z} U(\mathbf{r}'')$ (equation \ref{eq:potential-U-1st-deriv-space}) 
	on $\mathcal{S}''$.
	
	% Geophysical implications
	\section{Geophysical implications}
	
	In the present section, we discuss the implications and relations of our 
	theoretical results to well-established empirical approaches commonly found 
	in the geophysical literature about the equivalent-layer technique.
	
	\subsection{Joint processing of gravity-gradient data}
	
	The common approach to jointly process gravity-gradient data 
	\cite[e.g.,][]{boggs-dransfield2004,barnes_lumley2011, martinez_li_2016} or to convert gravity disturbance data into the 
	gravity-gradient tensor \cite[e.g.,][]{piauilino_etal2019_seg, oliveirajr_etal2023}, for example, 
	usually presume, without any proof, the existence of a single layer of 
	monopoles that satisfactorily reproduces the gravitational field and
	preserves the physical relationship with the true sources.
	This implies that the gravitational potential and its spatial derivatives attributed to the layer are presumed to 
	accurately reproduce the corresponding field components produced by the true sources.
	We provide the theoretical background for these potential-field transformations.
	
	Let us consider, for example, the common two-step approach to convert gravity disturbance into gravity
	gradient data (Table \ref{tbl:Table2}). The first step consists in obtaining a particular equivalent layer $p_{\varepsilon}$ 
	by considering that the observed potential field $d^{o}$ is given by $\tfrac{\partial}{\partial z} U(\mathbf{r})$ 
	(equation \ref{eq:potential-U-1st-deriv-space}) and that the kernel function 
	$a(\mathbf{r}, \mathbf{r}'')$ is the vertical derivative of $\Psi$ (equation \ref{eq:Psi}).
	The second step consists in using $p_{\varepsilon}$ to compute the gradient components.
	The $uv$-component of the gradient tensor is computed by setting $D_{f} \equiv \tfrac{\partial^{2}}{\partial u \, \partial v}$
	(equations \ref{eq:potential-U-2nd-deriv-space} and \ref{eq:transformed-field-continuous-layer-Neumann}),
	where $u, v \in \{ x, y, z \}$.
	
	Note that the first step consists in defining an equivalent layer 
	by setting $\alpha_{d} = \alpha_{a} \equiv \alpha(1, \hat{\mathbf{z}})$ in equation \ref{eq:p-Fourier}.
	According to Definition \ref{def:eql-Neumann}, this particular $p_{\varepsilon}$ is a Neumann's equivalent layer,
	which naturally preserves the physical relationship with the true sources.
	The second step of the common approach corresponds to 
	compute transformed fields $f_{\varepsilon}$ (equation \ref{eq:transformed-field-continuous-layer-Fourier}) 
	defined in the Fourier domain by using $\alpha_{t} = \alpha_{f}$ (equation \ref{eq:alpha-t-generic}), with 
	$\alpha_{f} \equiv \alpha(1, \hat{\mathbf{u}}) \, \alpha(1, \hat{\mathbf{v}})$ (equation \ref{eq:diff-operators-generic}), 
	where the unit vectors $\hat{\mathbf{u}}$ and $\hat{\mathbf{v}}$ are given by the canonical base vectors (equation \ref{eq:canonical-base-vector}).
	We can easily get from \ref{eq:diff-operators-generic} that this Fourier-domain operator $\alpha_{f}$ is equivalent to the 
	space-domain operator $D_{f} \equiv \tfrac{\partial^{2}}{\partial u \, \partial v}$.
	Thus, the common two-step approach to convert gravity disturbance into gravity
	gradient data via equivalent-layer technique can be theoretically explained as an application of 
	a particular Neumann's equivalent layer (Definition \ref{def:eql-Neumann}).
	
	
	\subsection{The existence of a planar dipole layer}
	
	
	A predominant trend in previous studies utilizing a planar equivalent layer of dipoles for magnetic-field transformations 
	is to treat the layer as a purely mathematical abstraction employed for processing total-field anomaly data (Table \ref{tbl:Table2}), 
	a perspective exemplified by the remarks of \citet{pedersen1989} concerning the research of 
	\citet{bhattacharyya_chan1977}.
	Specifically, Pedersen pointed out that the first fundamental existence theorem for the Dirichlet's 
	problem does not postulate that a general solution can be written as a layer of dipoles magnetized along its 
	normal direction. In essence, Pedersen underscored that a dipole layer lacks physical significance.
	In response, Chan affirmed that such a layer was employed in his work 
	solely to address an integral equation, i.e., as a purely mathematical solution.
	
	Since its inception in the 1960s, the equivalent-layer technique in applied geophysics has garnered 
	relatively scant attention in terms of exploring the theoretical attributes of planar dipole layers 
	capable of faithfully reproducing total-field anomaly data. 
	This inquiry usually aims to discern the physical significance and possible link of such layers 
	with boundary problems of potential theory.
	
	Here, we leverage a theoretical approach to infer within the purview of applied geophysics, 
	for the first time, the existence of planar dipole layers capable of accurately and simultaneously reproducing 
	the magnetic scalar potential and induction field components along any direction, including that representing 
	the approximated total-field anomaly (Table \ref{tbl:Table2}).
	
	We first consider that (i) the observed potential field $d^{o}$ 
	is satisfactorily described by the directional derivative $\hat{\mathbf{u}}^{\top} \mathbf{H}_{U}(\mathbf{r}) \, \hat{\mathbf{v}}$ 
	(equation \ref{eq:potential-U-oblique-2nd-deriv-space}), where $\hat{\mathbf{u}}$ defines the constant direction of the main geomagnetic
	field at the study area and $\hat{\mathbf{v}}$ the constant total-magnetization direction of the sources  (Table \ref{tbl:Table2}), and (ii)
	the kernel is $a(\mathbf{r}, \mathbf{r}'') = \hat{\mathbf{u}}^{\top} \mathbf{H}_{\Psi}(\mathbf{r}, \mathbf{r}'') \, \hat{\mathbf{h}}$,
	where $\mathbf{H}_{\Psi}$ is the Hessian matrix \citep[e.g.,][ p. 158]{kaplan2002} of
	function $\Psi$ (equation \ref{eq:Psi}) and $\hat{\mathbf{h}}$ is a constant unit vector not necessarily equal to 
	$\hat{\mathbf{v}}$.
	In this case, we use $\alpha_{d} \equiv \alpha(1, \hat{\mathbf{u}}) \, \alpha(1, \hat{\mathbf{v}})$ and 
	$\alpha_{a} \equiv \alpha(1, \hat{\mathbf{u}}) \, \alpha(1, \hat{\mathbf{h}})$ 
	(equation \ref{eq:diff-operators-generic}) in equation \ref{eq:p-Fourier} to obtain the following equivalent layer 
	in the Fourier domain:
	\begin{equation}
		\mathcal{F} [ p_{\varepsilon}(\mathbf{r}'') ] =
		k \, \tfrac{\left(\mathbf{k}^{\top} \hat{\mathbf{v}} \right)}{\left(\mathbf{k}^{\top} \hat{\mathbf{h}} \right)} \,
		\mathcal{F} \left[ U_{\varepsilon}(\mathbf{r}'') \right] \: .
		\label{eq:p-Fourier-dipoles}
	\end{equation}
	This particular equivalent layer represents a planar layer of dipoles with constant magnetization direction
	$\hat{\mathbf{h}}$ and was deduced by considering
	that $\hat{\mathbf{u}}$ is a unit vector having the same constant direction of the main geomagnetic field at the study area.
	However, it can be easily verified that the same equivalent layer can also be deduced by considering that $\hat{\mathbf{u}}$ is one of the
	canonical base vectors $\hat{\mathbf{x}}$, $\hat{\mathbf{y}}$ or $\hat{\mathbf{z}}$ (equation \ref{eq:canonical-base-vector}).
	In this case, the observed potential field represents a Cartesian component of the magnetic induction field  (Table \ref{tbl:Table2}).
	It can also be easily verified that the same $p_{\varepsilon}$ given by equation 
	\ref{eq:p-Fourier-dipoles} can be equally deduced if we consider that the 
	observed potential field is the magnetic scalar potential  (Table \ref{tbl:Table2}). 
	In this case, $\alpha_{d}$ and $\alpha_{a}$ are defined
	without the term $\alpha(1, \hat{\mathbf{u}})$. 
	
	The main conclusion obtained from the reasoning presented above is that a planar equivalent layer of dipoles defined according 
	to equation \ref{eq:p-Fourier-dipoles} simultaneously 
	reproduces the magnetic scalar potential and any magnetic induction field component (including the approximated total-field anomaly) 
	due to sources with constant total-magnetization direction defined by $\hat{\mathbf{v}}$, even if the layer has a constant magnetization 
	direction defined by a unit vector $\hat{\mathbf{h}} \ne \hat{\mathbf{v}}$. 
	This theoretical result explains, for example, why we can use a planar equivalent layer of dipoles with arbitrary constant 
	magnetization direction $\hat{\mathbf{h}}$ to compute the magnetic induction field components or the amplitude of the magnetic induction field due to the true sources 
	without presuming any information about the constant direction of their total magnetization 
	\citep[e.g,][]{li_etal_2021}.
	
	
	\subsection{Computing the RTP by using a dipole layer}
	
	
	According to the pioneer work of \citet{silva1986}, which is widely followed in applied geophysics, 
	the RTP in the space domain is calculated from total-field anomaly data (Table \ref{tbl:Table2}) in two steps.
	The first consists in estimating the magnetic-moment intensity distribution on a planar layer 
	of dipoles (or doublets) with constant magnetization direction equal to that of the true sources.
	The second step consists in computing the $z$-component of the magnetic induction field due 
	to the estimated layer by setting its magnetization along the vertical direction.
	It is usually verified a posteriori that this computed field approximates the true RTP anomaly (Table \ref{tbl:Table2}).
	
	We call attention to the fact that \citet{silva1986} cited the work of \citet{roy1962} to justify that ``the magnetic 
	field due to arbitrary sources can be exactly fitted by the magnetic field due to a horizontal equivalent
	layer of dipoles''. 
	However, \citeauthor{roy1962} neither discussed the existence of a dipole layer that fits the observed total-field 
	anomaly data nor its use to retrieve the RTP anomaly, but 
	presented a general outline to deduce the solutions of the Dirichlet's and Neumann's problems on a plane 
	located above the sources (equation \ref{eq:U-Dirichlet-Neumann}).
	In other words, the two-step approach proposed by \citet{silva1986}, which is widely accepted by the scientific community 
	and we know to be correct from practice 
	\citep[e.g., ][]{leao-silva_1989, guspi_novara2009, oliveirajr_etal2013, li_etal2014, reis_etal2020, takahashi-etal2022, piauilino_etal2025}, 
	is an empirical generalization of the \citeauthor{roy1962}'s (\citeyear{roy1962}) work. 
	It is also important to point out that the mathematical equivalence between the space-domain approach to compute the 
	RTP via equivalent-layer technique and the classical approach directly using the Fourier transform of the observed 
	total-field anomaly data \citep[e.g., ][ p. 331]{blakely1996} have never been clearly established in the literature.
	Here, we provide a general approach aiming at filling these theoretical gaps.
	
	From a theoretical point of view, the two-step approach to compute the RTP in the space domain 
	consists in first setting a planar layer of dipoles according to equation \ref{eq:p-Fourier-dipoles}, 
	by considering the unit vector $\hat{\mathbf{h}}$ equal to $\hat{\mathbf{v}}$, which defines 
	the constant total-magnetization direction of the true sources.
	In this case, the ratio $\nicefrac{\left(\mathbf{k}^{\top} \hat{\mathbf{v}} \right)}{\left(\mathbf{k}^{\top} \hat{\mathbf{h}} \right)} = 1$ 
	and we can easily verify that the $\mathcal{F} [ p_{\varepsilon}(\mathbf{r}'') ]$ given by 
	equation \ref{eq:p-Fourier-dipoles} assumes the form of a Neumann's equivalent layer (Definition \ref{def:eql-Neumann}).
	The second step to compute the RTP consists in using the predefined Neumann's layer to compute 
	the transformed potential field $f_{\varepsilon}$ (equation \ref{eq:transformed-field-continuous-layer}) defined in the Fourier domain 
	by using $\alpha_{t} = \alpha_{f}$ (equation \ref{eq:alpha-t-generic}), with 
	$\alpha_{f} \equiv \alpha(2, \hat{\mathbf{z}})$ (equations \ref{eq:diff-operators-generic} and \ref{eq:canonical-base-vector}).
	In this case, it can be verified from equations \ref{eq:diff-operators-generic} and \ref{eq:transformed-field-continuous-layer-single-step} 
	that the computed transformed field $f_{\varepsilon}$ (equation \ref{eq:transformed-field-continuous-layer}) 
	approximates the RTP anomaly  (Table \ref{tbl:Table2}). 
	As we have just shown here, for the first time, the classical approach of \citet{silva1986} to compute the RTP in the 
	space domain can be formally explained in terms of a Neumann's equivalent layer (Definition \ref{def:eql-Neumann}). 
	
	We can also investigate the theoretical link between the RTP in the space domain, 
	via equivalent-layer technique, and in the Fourier domain, via classical approach \citep[e.g., ][ p. 331]{blakely1996}. 
	As we have already discussed, the transformed field $f_{\varepsilon}$ computed via \ref{eq:transformed-field-continuous-layer-Fourier} 
	from a previously defined equivalent layer $p_{\varepsilon}$ can also be directly 
	computed from the observed potential field $d^{o}$ (equation \ref{eq:transformed-field-continuous-layer-1-single-step}
	or \ref{eq:transformed-field-continuous-layer-2-single-step}), 
	in terms of the ratio $\nicefrac{\alpha_{f}}{\alpha_{d}}$.
	To compute the RTP, we set $\alpha_{d} \equiv \alpha(1, \hat{\mathbf{u}}) \, \alpha(1, \hat{\mathbf{v}})$ 
	and $\alpha_{f} \equiv \alpha(2, \hat{\mathbf{z}})$ (equation \ref{eq:diff-operators-generic}).
	By using these specific $\alpha_{d}$ and $\alpha_{f}$, setting $z^{\ddagger}_{0} = z^{\dagger}_{0}$ and 
	(i) letting $k_{\varepsilon}$ goes to infinity in equation \ref{eq:transformed-field-continuous-layer-1-single-step} or
	(ii) letting $\mu_{\varepsilon}$ goes to zero in equation \ref{eq:transformed-field-continuous-layer-2-single-step}, we get 
	\begin{equation}
		\mathcal{F} [ f_{\varepsilon}(\mathbf{r}^{\dagger}) ] =
		\tfrac{k^{2}}{\left(\mathbf{k}^{\top} \hat{\mathbf{u}} \right) \, \left(\mathbf{k}^{\top} \hat{\mathbf{v}} \right)} \, 
		\mathcal{F} \left[ d^{o}(\mathbf{r}^{\dagger}) \right] \: ,
		\label{eq:transformed-field-continuous-layer-single-step-RTP}
	\end{equation}
	which represents the
	classical approach to compute the RTP in the Fourier domain \citep[e.g., ][ p. 331]{blakely1996}.
	We stress two important aspects of this result. The first concerns the theoretical link 
	between the space- and Fourier-domain approaches to compute the RTP. To the authors knowledge, this link has never been fully 
	detailed in the literature.
	The second important aspect of our theoretical approach is that it provides the 
	reason by which the equivalent-layer technique ultimately suffers the same instability as 
	the Fourier-domain RTP operator at low latitudes. In this case, the estimated equivalent layer of dipoles 
	presents striations, as have already been pointed out by \citet{li_etal2014}. 
	In the Fourier domain, these artifacts are directly related to the term 
	$\nicefrac{\left(\mathbf{k}^{\top} \hat{\mathbf{v}} \right)}{\left(\mathbf{k}^{\top} \hat{\mathbf{h}} \right)}$ 
	in equation \ref{eq:p-Fourier-dipoles}. This term produces, in the equivalent layer of dipoles, a 
	numerical instability similar to that caused by the term 
	$\nicefrac{k^{2}}{\left(\mathbf{k}^{\top} \hat{\mathbf{u}} \right) \, \left(\mathbf{k}^{\top} \hat{\mathbf{v}} \right)}$ in equation 
	\ref{eq:transformed-field-continuous-layer-single-step-RTP},
	which is the base of the classical approach to compute the RTP in the Fourier domain \citep[e.g., ][ p. 331]{blakely1996}.
	
	
	\subsection{Generalized equivalent sources}
	
	Let us define an equivalent layer $p_{\varepsilon}$ by considering the kernel function 
	$a(\mathbf{r}, \mathbf{r}'') = \Psi(\mathbf{r}, \mathbf{r}'')$ (equation \ref{eq:Psi}). 
	Note that, in this case, 
	the corresponding predicted potential field $d_{\varepsilon}$ (equation \ref{eq:predicted-data-continuous-layer}) 
	is analogous to the gravitational potential due to the equivalent layer.
	Besides, according to equations \ref{eq:p-Fourier}--\ref{eq:kernel-a-Fourier}, we have that 
	$\alpha_{a} = 1$ and 
	\begin{equation}
		\mathcal{F} [ p_{\varepsilon}(\mathbf{r}'') ] =
		k \, \alpha_{d} \, \mathcal{F} \left[ U_{\varepsilon}(\mathbf{r}'') \right]  \: ,
		\label{eq:p-Fourier-generalized-eqs}
	\end{equation}
	which approximates the vertical derivative of the observed potential field on the layer plane. 
	Such layers produce potential-field transformations (equations \ref{eq:transformed-field-continuous-layer} and 
	\ref{eq:transformed-field-continuous-layer-Fourier}) defined by $\alpha_{t} = \nicefrac{\alpha_{f}}{\alpha_{d}}$ (equation \ref{eq:alpha-t-generic}).
	
	This strategy, called \textit{generalized equivalent sources}, was proposed by \citet{cordell1992} and recently adapted by 
	\citet{soler_uieda2021} to perform gravity-data processing. 
	Those authors considered only the particular case in which the observed potential field $d^{o}$ is gravity disturbance data 
	$\tfrac{\partial}{\partial z} U(\mathbf{r})$  (Table \ref{tbl:Table2}), 
	with $\alpha_{d} \equiv \alpha(1, \hat{\mathbf{z}}) = k$ (equations \ref{eq:observed-data-Fourier}, 
	\ref{eq:diff-operators-generic} and \ref{eq:canonical-base-vector}).
	In this particular case, equation \ref{eq:p-Fourier-generalized-eqs} leads to an equivalent layer that 
	(i) approximates $\tfrac{\partial^{2}}{\partial z^{2}} U(\mathbf{r})$ (equation \ref{eq:potential-U-2nd-deriv-space}) 
	in the space domain, 
	(ii) has minimum phase changes (Definition \ref{def:eql-min-phase}) and 
	(iii) results in potential-field transformations defined by $\alpha_{t} = \nicefrac{\alpha_{f}}{k}$ (equation \ref{eq:alpha-t-generic}).
	
	\subsection{Alternative approach to compute the RTP}
	
	\citet{guspi_novara2009} proposed an alternative two-step approach to compute the RTP via equivalent-layer technique. 
	The first step consists in obtaining an equivalent layer by considering that (i) the observed potential field $d^{o}$
	is satisfactorily described by the directional derivative 
	$\hat{\mathbf{u}}^{\top} \mathbf{H}_{U}(\mathbf{r}) \, \hat{\mathbf{v}}$ 
	(equation \ref{eq:potential-U-oblique-2nd-deriv-space}), where $\hat{\mathbf{u}}$ defines the constant direction of the main geomagnetic
	field at the study area and $\hat{\mathbf{v}}$ the constant total-magnetization direction of the sources  (Table \ref{tbl:Table2}), and (ii)
	the kernel function is given by 
	\begin{equation}
		a(\mathbf{r}, \mathbf{r}'') = \tfrac{\left( z''_{0} - z \right)}{2 \, \pi} \, 
		\ln{\Big[} \left( z''_{0} - z \right) + \| \mathbf{r}'' - \mathbf{r} \| {\Big]} - 
		\tfrac{\| \mathbf{r}'' - \mathbf{r} \|}{2 \, \pi} \: ,
		\label{eq:kernel-Fourier-Guspi}
	\end{equation}
	which represents the double integral of the function $\Psi$ (equation \ref{eq:Psi}) 
	along the vertical direction $z$.
	The second step of the alternative approach proposed by \citet{guspi_novara2009} computes the RTP anomaly 
	(Table \ref{tbl:Table2}) by setting $f_{\varepsilon}$ (equation \ref{eq:transformed-field-continuous-layer}) with 
	$t(\mathbf{r}, \mathbf{r}'') = \Psi(\mathbf{r}, \mathbf{r}'')$ (equation \ref{eq:Psi}).
	
	Note that the first step consists in obtaining an equivalent layer $p_{\varepsilon}$ defined by 
	$\alpha_{d} = (\mathbf{k}^{\top} \hat{\mathbf{u}}) (\mathbf{k}^{\top} \hat{\mathbf{v}})$ (equations \ref{eq:observed-data-Fourier} 
	and \ref{eq:diff-operators-generic}) and 
	$\alpha_{a} = \nicefrac{(\mathbf{k}^{\top} \hat{\mathbf{u}}) (\mathbf{k}^{\top} \hat{\mathbf{h}})}{k^{2}}$, 
	where the term $\nicefrac{1}{k^{2}}$ represents the double integration along the vertical direction. In this case,  
	$p_{\varepsilon}$ assumes the following form in the Fourier domain:
	\begin{equation}
		\mathcal{F} [ p_{\varepsilon}(\mathbf{r}'') ] =
		k^{3} \, \tfrac{\left(\mathbf{k}^{\top} \hat{\mathbf{v}} \right)}{\left(\mathbf{k}^{\top} \hat{\mathbf{h}} \right)} \,
		\mathcal{F} \left[ U_{\varepsilon}(\mathbf{r}'') \right] \: .
		\label{eq:p-Fourier-guspi}
	\end{equation}
	The second step consists in computing $f_{\varepsilon}$ in the Fourier domain (equation \ref{eq:transformed-field-continuous-layer-Fourier}) 
	by using $\alpha_{t} = 1$ (equation \ref{eq:alpha-t-generic}).
	In this case, it can be verified that the transformed potential field $\mathcal{F} [ f_{\varepsilon}(\mathbf{r}^{\ddagger}) ]$ 
	(equation \ref{eq:transformed-field-continuous-layer-Fourier}) actually approximates the RTP anomaly  (Table \ref{tbl:Table2}) 
	for the particular case in which $\hat{\mathbf{h}} = \hat{\mathbf{v}}$, 
	i.e., for the case in which the constant magnetization direction of the equivalent layer coincides with that of the true sources.
	This result shows the effectiveness of our approach in providing the theoretical background not only for classical 
	equivalent-layer methods, but also for unconventional ones.
	
	% Numerical simulations
	\section{Numerical simulations}
	\label{eq:numerical_simulations}
	
	\subsection{General setup}
	
	We present three numerical tests illustrating planar equivalent layers that intersect a simulated 
	magnetic source at the same horizontal plane.
	Each test involves defining three versions of a predefined equivalent layer by following three different approaches.
	The first approach involves estimating a discrete equivalent layer by solving a linear inverse problem. 
	This is the common strategy found in literature.
	The second and third approaches estimate discrete layers by numerically evaluating equation \ref{eq:p-Fourier}, 
	as we explain later in the present section. 
	It is important to bear in mind that the first approach is the only one that can be applied to field data.
	The equivalent layers obtained by following the second and third approaches depend on the knowledge of
	the position, shape and physical-property distribution of the true sources. Consequently, they cannot be 
	obtained from field data produced by unknown sources, but only from numerical simulations.
	Our objective is to demonstrate that the theoretical framework established here for continuous equivalent 
	layers is consistent with the classical space-domain approach, which involves solving a linear inverse problem.
	Here we define the general setup for all tests.
	
	\subsubsection{Synthetic magnetic source}
	
	We have simulated a magnetic source composed of $1241$ horizontally juxtaposed rectangular prisms
	(Figure \ref{fig:Figure2}). 
	The prisms extend along the $x$-direction from $\approx -3601$ to $4900$ m and 
	along the $y$-direction from $\approx -3000.3$ to $\approx 5800.2$ m. 
	Each prism has horizontal side lengths of approximately $256.4$ m in the $x$-direction 
	and $204.1$ m in the $y$-direction, with thicknesses ranging from $0$ m to $600$ m. 
	
	The top of the source is delimited by the plane $z^{t}_{0} = 0$ m and varies 
	from $0$ to $300$ m, while its bottom ranges from $300$ to $600$ m.
	The source has a pseudo-random total-magnetization intensity $\sigma(\mathbf{r}')$ (Table \ref{tbl:Table1}) 
	varying from $2$ to $3$ A/m according to a uniform distribution and a constant direction defined by the 
	unit vector $\hat{\mathbf{v}}$, with inclination $-10.2^{\circ}$ and declination $-19.8^{\circ}$. 
	We also consider that the main geomagnetic field has a constant direction given by the unit vector $\hat{\mathbf{u}}$,
	with inclination $-30^{\circ}$ and declination $-19^{\circ}$. 
	
	\subsubsection{Synthetic potential fields}
	
	We have set a regular grid of $N_{x} \times N_{y}$ observation points 
	$\mathbf{r}^{\dagger}_{i} = (x^{\dagger}_{i}, y^{\dagger}_{i}, z^{\dagger}_{0})^{\top}$, $i \in \{1, \dots, N = N_{x} \, N_{y}\}$,
	where $N_{x} = 400$, $N_{y} = 330$. 
	The points are located on a planar surface $\mathcal{S}^{\dagger}$ (Figure \ref{fig:Figure1}) above the magnetic source
	(Figure \ref{fig:Figure2}), at $z^{\dagger}_{0} = -150$ m.
	Figure \ref{fig:Figure3} shows the potential fields (Table \ref{tbl:Table2}) produced by the magnetic source at 
	the simulated grid of observation points.
	The potential fields were computed by using the closed expressions derived by \cite{Nagyetal2000}, \cite{Nagyetal2002} 
	and \cite{Fukushima2020}.
	
	In our three tests, we consider that the observed potential field $d^{o}$ (equations \ref{eq:observed-data-space} and \ref{eq:observed-data-Fourier}) 
	is the approximated total-field anomaly (Figure \ref{fig:Figure3}a), which means that 
	$\alpha_{d} \equiv \alpha(1, \hat{\mathbf{u}}) \, \alpha(1, \hat{\mathbf{v}})$ (equation \ref{eq:diff-operators-generic}).
	
	\subsubsection{Equivalent layers layout}
	
	In all tests, we define discrete equivalent layers on the same planar surface $\mathcal{S}''$ (Figure \ref{fig:Figure1}), 
	that intersects the source (Figure \ref{fig:Figure2}) at $z''_{0} = 200$ m.
	All layers are composed of punctual equivalent sources arranged on a regular grid of $N_{x} \times N_{y}$ points 
	$\mathbf{r}''_{j} = (x''_{j}, y''_{j}, z''_{0})^{\top}$, $j \in \{1, \dots, M = N_{x} \, N_{y}\}$,
	so that each point $\mathbf{r}''_{j}$ is located right below an observation point $\mathbf{r}^{\dagger}_{i}$.
	
	\subsection{Tests}
	
	Each test consists in first specifying the equivalent layer type by setting the operators $\alpha_{a}$ and $D_{a}$ 
	via equation \ref{eq:diff-operators-generic}.
	They act on the scaled inverse distance function $\Psi$ 
	(equations \ref{eq:Psi} and \ref{eq:Psi-Fourier}) in the space and Fourier domains,
	respectively, to define the elements $a_{ij}$ of matrix $\mathbf{A}$ (Definition \ref{def:discrete-eqlayer}) and 
	the kernel function $a(\mathbf{r}, \mathbf{r}'')$ (Definition \ref{def:continuous-eqlayer}).
	The second step of each test is determining three versions of the predefined equivalent layer 
	on the planar surface $\mathcal{S}''$ by following three different approaches.
	
	\subsubsection{Approach 1: Classical space-domain equivalent layer}
	
	The first approach involves solving a linear inverse problem to estimate a discrete 
	equivalent layer $\mathbf{p}_{\epsilon}$ leading to a predicted data $\mathbf{d}_{\epsilon}$ 
	close to the observed data $\mathbf{d}^{o}$ (Definition \ref{def:discrete-eqlayer}).
	In this case, the discrete layer is totally defined in the space domain, by following the 
	classical approach found in the literature.
	
	To solve the linear inverse problem, we apply the method proposed by 
	\citet{takahashi-etal2020, takahashi-etal2022} because the spatial configuration of the 
	simulated equivalent sources and observation points leads to 
	matrices $\mathbf{A}$ (Definition \ref{def:discrete-eqlayer}) having a Block-Toeplitz 
	Toeplitz-Block (BTTB) structure.
	This method estimates $\mathbf{p}_{\epsilon}$ by using the Conjugate Gradient Least Squares 
	(CGLS) algorithm, which naturally imposes some degree of smoothness when run few iterations 
	\citep[][ p. 166]{aster_etal2019}.
	In fact, the discrete layers estimated here were obtained with only $100$ iterations of the 
	CGLS method, which is very small compared to the total number of equivalent sources 
	($N = 132,000$).
	
	\subsubsection{Approaches 2 and 3: Theoretical closed-form equivalent layers}
	
	The second and third approaches to define a given equivalent layer are developed in the Fourier domain 
	and bypass the need to solve a linear inverse problem. Instead, they directly approximate the continuous 
	equivalent layer $p_{\varepsilon}$ defined by equation \ref{eq:p-Fourier}. 
	
	In both approaches, we compute discrete layers according to the following steps.
	We start by computing the pseudogravity anomaly (Table \ref{tbl:Table2}) due to the synthetic source 
	at the simulated grid of observation points $\mathbf{r}^{\dagger}_{i}$ located above the synthetic source (Figure \ref{fig:Figure2}), 
	on the planar surface $\mathcal{S}^{\dagger}$ at $z^{\dagger}_{0} = -150$ m.
	The pseudogravity anomaly is defined by $k \, \mathcal{F}\left[ U(\mathbf{r}^{\dagger}) \right]$ in the Fourier domain,
	where $\mathcal{F}\left[ U(\mathbf{r}^{\dagger}) \right]$ is given by equation \ref{eq:potential-U-Fourier}.
	Then we numerically compute the term $k \, \mathcal{F}\left[ U_{\varepsilon}(\mathbf{r}'') \right]$ (equation \ref{eq:p-Fourier}).
	The difference between approaches 2 and 3 is the way we numerically compute $\mathcal{F}\left[ U_{\varepsilon}(\mathbf{r}'') \right]$.
	In approach 2, we use the equation 
	\ref{eq:potential-U-epsilon-Fourier-layer-1}, with a given truncating wavenumber $k_{\varepsilon}$. 
	In approach 3, the term $\mathcal{F}\left[ U_{\varepsilon}(\mathbf{r}'') \right]$ is computed by using equation 
	\ref{eq:potential-U-epsilon-Fourier-layer-2}, with a given regularization parameter $\mu_{\varepsilon}$.
	In both approaches, we obtain a discrete equivalent layer $\mathbf{p}_{\epsilon}$ (Definition \ref{def:discrete-eqlayer})
	in the space domain by numerically evaluating the inverse Fourier transform of equation \ref{eq:p-Fourier}, 
	using the predefined operators $\alpha_{d}$ and $\alpha_{a}$.
	
	Because the planar equivalent-layer theorem (Theorem \ref{theorem:planar-eqlayer-theorem}) is a nonconstructive existence statement 
	that does not define how to obtain $k_{\varepsilon}$ or $\mu_{\varepsilon}$ (equations \ref{eq:potential-U-epsilon-Fourier-layer-1} and 
	\ref{eq:potential-U-epsilon-Fourier-layer-2}), we adopt an empirical procedure in our numerical simulations.
	In approaches 2 and 3, we actually obtain sets of tentative discrete equivalent layers $\mathbf{p}_{\epsilon}$ (Definition \ref{def:discrete-eqlayer}) 
	for different values of $k_{\varepsilon}$ and $\mu_{\varepsilon}$. 
	The optima $k_{\varepsilon}$ and $\mu_{\varepsilon}$ are chosen as those that have minimum values 
	and produce discrete equivalent layers $\mathbf{p}_{\epsilon}$ that satisfactorily fit the approximated 
	total-field anomaly (Figure \ref{fig:Figure3}a).
	
	\subsection{Results}
	
	\subsubsection{Estimated equivalent layers}
	
	Figure \ref{fig:Figure4} shows all estimated equivalent layers obtained in our three tests.
	The first, second and third columns of Figure \ref{fig:Figure4} contain the results obtained from 
	each numerical test, while the first, second and third rows of Figure \ref{fig:Figure4} show the 
	results obtained by following approaches 1--3, respectively.
	As we can see in Figure \ref{fig:Figure5}, the equivalent layers obtained in each test, by following 
	different approaches, are very similar to each other.
	The differences oscillate around zero and are dominated by short wavelengths.
	
	The layers obtained from our first test (Figures \ref{fig:Figure4}a, d and g) are Neumann's 
	equivalent layers (Definition \ref{def:eql-Neumann}), 
	were set with $\alpha_{a} = \alpha_{d}$ (equation \ref{eq:p-Fourier}) and approximate the
	pseudogravity anomaly $\nicefrac{\partial U(\mathbf{r}'')}{\partial z}$ 
	(equation \ref{eq:potential-U-oblique-1st-deriv-space}) 
	due to the synthetic source (Figure \ref{fig:Figure2}) on the layer plane $\mathcal{S}''$.
	As expected, the estimated Neumann's layers have all positive magnetic moment distributions 
	with maxima located right above the synthetic source.
	The Neumann's layers obtained by following approaches 2 and 3 
	(Figures \ref{fig:Figure4}d and g) were defined by using a truncating wavenumber 
	$k_{\varepsilon} = 0.020$ (equation \ref{eq:potential-U-epsilon-Fourier-layer-1})
	and a regularization parameter $\mu_{\varepsilon} = 5.854$ 
	(equation \ref{eq:potential-U-epsilon-Fourier-layer-2}). 
	The $k_{\varepsilon}$ and $\mu_{\varepsilon}$ values were chosen by using the curves shown 
	in Figures \ref{fig:Figure5}a and d.
	
	In our second test, we estimate discrete layers (shown in Figures \ref{fig:Figure4}b, e and h) 
	with $\alpha_{a} = \left(\mathbf{k}^{\top} \hat{\mathbf{u}} \right) k$ 
	(equation \ref{eq:p-Fourier}).
	From the physical point of view, such layers represent the negative magnetic scalar 
	potential (Table \ref{tbl:Table2}) due to the synthetic source on the plane $\mathcal{S}''$.
	These layers oscillate from negative to positive values and approximate 
	$\nabla U(\mathbf{r}'')^{\top} \hat{\mathbf{v}}$ 
	(equation \ref{eq:potential-U-oblique-1st-deriv-space}) on the layer plane. 
	The layers obtained in the second test, by following approaches 2 and 3 
	(Figures \ref{fig:Figure4}e and h), were defined by using a truncating wavenumber 
	$k_{\varepsilon} = 0.025$ (equation \ref{eq:potential-U-epsilon-Fourier-layer-1})
	and a regularization parameter $\mu_{\varepsilon} = 1.000$ 
	(equation \ref{eq:potential-U-epsilon-Fourier-layer-2}). 
	The $k_{\varepsilon}$ and $\mu_{\varepsilon}$ values were chosen by using the curves shown 
	in Figures \ref{fig:Figure5}b and e.
	
	Figures \ref{fig:Figure4}c, f and i show the discrete equivalent layers estimated in our 
	third numerical test, with $\alpha_{a} = \left(\mathbf{k}^{\top} \hat{\mathbf{u}} \right)$ 
	(equation \ref{eq:p-Fourier}) 
	These layers approximate the $z$ component 
	$\frac{\partial}{\partial z}\nabla U(\mathbf{r}'')^{\top} \hat{\mathbf{v}}$
	of the magnetic induction field (Table \ref{tbl:Table2}) due to the synthetic source on the 
	layer plane.
	Similarly to the layers obtained in the second test (Figures \ref{fig:Figure4}b, e and h), 
	the layers obtained in the third test (Figures \ref{fig:Figure4}c, f and i) oscillate from
	negative to positive values. 
	The main difference is that the latter exhibit significantly shorter wavelength patterns that
	are concentrated over the source boundary.
	The layers obtained in the third test, by following approaches 2 and 3 
	(Figures \ref{fig:Figure4}f and i), were defined by using a truncating wavenumber 
	$k_{\varepsilon} = 0.015$ (equation \ref{eq:potential-U-epsilon-Fourier-layer-1})
	and a regularization parameter $\mu_{\varepsilon} = 49.889$ 
	(equation \ref{eq:potential-U-epsilon-Fourier-layer-2}). 
	The $k_{\varepsilon}$ and $\mu_{\varepsilon}$ values were chosen by using the curves shown 
	in Figures \ref{fig:Figure5}c and f.
	
	\subsubsection{Transformed potential fields}
	
	We computed the transformed potential fields $\mathbf{f}_{\epsilon}$ (equation \ref{eq:transformed-field-discrete-layer}) 
	at the observation points $\mathbf{r}^{\dagger}_{i}$ located on the plane $z^{\dagger}_{0} = -150$ m by using the discrete 
	layers obtained from linear inversion (Figures \ref{fig:Figure4}a–c). Figure \ref{fig:Figure6} shows the residuals 
	between the potential fields generated by these estimated layers and the synthetic source. The transformed potential 
	fields generated by the remaining estimated layers (shown in Figures \ref{fig:Figure4}d–i) are similar to those in Figure 
	\ref{fig:Figure6} and are provided as supplementary material.
	
	To compute the approximated total-field anomaly and the $x$, $y$ and $z$ components of the magnetic induction field 
	due to the Neumann's layers (Figures \ref{fig:Figure4}a, d and g), 
	we have used matrices $\mathbf{T}$ (equation \ref{eq:transformed-field-discrete-layer}) 
	having elements $t_{ij} = \hat{\mathbf{t}}^{\top} \mathbf{H}_{\Psi}(\mathbf{r}^{\dagger}_{i}, \mathbf{r}''_{j}) \, \hat{\mathbf{h}}$,
	with unit vectors $\hat{\mathbf{t}}$ equal to, respectively, $\hat{\mathbf{u}}$, $\hat{\mathbf{x}}$, $\hat{\mathbf{y}}$ and
	$\hat{\mathbf{z}}$ (equation \ref{eq:canonical-base-vector}), 
	where $\mathbf{H}_{\Psi}$ is the Hessian of function $\Psi(\mathbf{r}^{\dagger}_{i}, \mathbf{r}''_{j})$
	(equation \ref{eq:Psi}).
	To compute the RTP anomaly, we have used a matrix $\mathbf{T}$ 
	having elements $t_{ij} = \tfrac{\partial^{2}}{\partial z^{2}} \Psi(\mathbf{r}^{\dagger}_{i}, \mathbf{r}''_{j})$.
	
	The approximated total-field anomaly and $x$, $y$ and $z$ components of the magnetic induction field 
	due to the layers shown in Figures \ref{fig:Figure4}b, e and h were computed 
	by following the same approach for the Neumann's layer. 
	The RTP, however, was computed differently, by following two steps. 
	The first step was defined in the space domain and involved multiplying the estimated discrete layer 
	by a matrix $\mathbf{T}$ 
	having elements $t_{ij} = \tfrac{\partial^{2}}{\partial z^{2}} \Psi(\mathbf{r}^{\dagger}_{i}, \mathbf{r}''_{j})$,
	which is the same approach used for the Neumann's layer. 
	The resulting transformed field was then numerically convolved with the filter $\nicefrac{k}{\left(\mathbf{k}^{\top} \hat{\mathbf{v}} \right)}$ 
	in the Fourier domain, where $\mathbf{v}$ defines the constant magnetization direction of the synthetic source.
	This second step is necessary to include the factor $\nicefrac{\alpha_{a}}{\alpha_{d}}$ into the operator 
	$\alpha_{t}$ associated with the transformed field (equation \ref{eq:alpha-t-generic}).
	
	The estimated discrete layers shown in Figure \ref{fig:Figure4}c, f and i do not represent the conventional 
	equivalent-layer approach applied to magnetic data processing.
	Because of that, they lead to potential-field transformations that are not commonly 
	found in the literature about equivalent-layer methods.
	To compute the approximated total-field anomaly, we use a matrix $\mathbf{T}$ (equation \ref{eq:transformed-field-discrete-layer}) 
	with elements $t_{ij} = \nabla \Psi(\mathbf{r}^{\dagger}_{i}, \mathbf{r}''_{j})^{\top} \hat{\mathbf{u}}$, 
	which represent the directional derivative of first order of $\Psi$ along the constant direction $\hat{\mathbf{u}}$ 
	of the simulated main field.
	The $x$, $y$ and $z$ components of the magnetic induction field were computed by using matrices $\mathbf{T}$
	having elements $t_{ij} = \tfrac{\partial}{\partial x} \Psi(\mathbf{r}^{\dagger}_{i}, \mathbf{r}''_{j})$, 
	$t_{ij} = \tfrac{\partial}{\partial y} \Psi(\mathbf{r}^{\dagger}_{i}, \mathbf{r}''_{j})$ and 
	$t_{ij} = \tfrac{\partial}{\partial z} \Psi(\mathbf{r}^{\dagger}_{i}, \mathbf{r}''_{j})$, respectively.
	To compute the RTP anomaly due to the estimated layers shown in Figures \ref{fig:Figure4}c, f and i, we performed two steps. 
	The first step was defined in the space domain and involved multiplying the estimated discrete layer 
	by a matrix $\mathbf{T}$ 
	having elements $t_{ij} = \tfrac{\partial^{2}}{\partial z^{2}} \Psi(\mathbf{r}^{\dagger}_{i}, \mathbf{r}''_{j})$,
	which is the same approach used for the Neumann's layer. 
	The resulting transformed field was then numerically convolved with the filter $\nicefrac{1}{\left(\mathbf{k}^{\top} \hat{\mathbf{v}} \right)}$ 
	in the Fourier domain, where $\mathbf{v}$ defines the constant magnetization direction of the synthetic source.
	This second step is necessary to include the factor $\nicefrac{\alpha_{a}}{\alpha_{d}}$ into the operator 
	$\alpha_{t}$ associated with the transformed field (equation \ref{eq:alpha-t-generic}).
	
	\subsubsection{Remarks on the numerical results}
	
	There are three key aspects of our results that we would like to highlight.
	First, the differences shown in Figure \ref{fig:Figure5} demonstrate that our theoretical 
	framework is compatible with the classical approach for estimating equivalent layers in the space 
	domain through the solution of a linear inverse problem.
	In particular, the estimated layers obtained from approaches 2 and 3  -- 
	by numerically evaluating the inverse Fourier transform of equation \ref{eq:p-Fourier} --
	are closely aligned with those derived from approach 1, which involves solving a linear inverse problem.
	
	Another important aspect of our results is that all estimated layers (Figure \ref{fig:Figure4}) 
	produce transformed potential fields $\mathbf{f}_{\epsilon}$ (equation \ref{eq:transformed-field-discrete-layer})  
	that closely match the fields generated by the synthetic source (Figure \ref{fig:Figure2}), 
	as shown in Figure \ref{fig:Figure6}. The residuals for the RTP anomalies 
	(Figures \ref{fig:Figure6}g, h, and i) exhibit trends along the constant magnetization direction of 
	the synthetic source. These trends arise from the relatively low inclination of the synthetic source 
	magnetization, $-10.2^{\circ}$.
	
	Finally, our results demonstrate that potential-field transformations can be computed with any 
	equivalent layer, provided that $\alpha_{t}$ (equation \ref{eq:alpha-t-generic}) is appropriately 
	set. This means that the reduction to the pole, for example, is not a transformation that can only 
	be computed using the equivalent-layer technique with layers formed by dipoles. Instead, our 
	results show that other layers can also be used to estimate the RTP anomaly, as long as $\alpha_{t}$ 
	(equation \ref{eq:alpha-t-generic}) is properly adjusted, thus confirming our theoretical findings. 
	We have shown that when the layer is not formed by dipoles, information about the constant directions 
	of the main field and source magnetization is incorporated into $\alpha_{t}$ 
	(equation \ref{eq:alpha-t-generic}), which the interpreter uses to compute the desired 
	potential-field transformation.
	
	% Application to field data
	\section{Application to field data}
	\label{eq:field_data}
	
	We have applied the equivalent-layer technique to the Registro do Araguaia anomaly 
	\citep{dutra_etal_2012, marangoni_mantovani_2013, dutra_etal_2014}, which is one the most 
	prominent anomalies over the NW portion of the Goiás Alkaline Province (GAP), Brazil 
	\citep{junqueira-brod_etal_2005}. 
	According to \citet{dutra_etal_2012}, there are no outcrops or any geological feature 
	indicating the source of this anomaly and the region is totally covered by quaternary 
	sediments from the Araguaia River.
	The GAP is a result of mafic-alkalic magmatism that occurred in the Late Cretaceous. Their 
	intrusions penetrate the Bras{\'i}lia Belt basement and the Goi{\'a}s Magmatic Arch, as well 
	as intrude and overlie basalts and sedimentary rocks from the northern border of the 
	Paran{\'a} Basin.
	\citet{dutra_etal_2014} consider that the source of Registro do Araguaia anomaly is one of the 
	most important magma feeders for the GAP intrusions. 
	They propose that this anomaly and its NE-SW lineament may be the main locus for magma conduit 
	at the southern end of the Transbrasiliano Lineament.
	In fact, this lineament and the Bom Jardim de Goiás Arch are recognized as the major 
	Precambrian structures related with the ascent of the Cretaceous alkaline magmas of GAP 
	\citep[][ p. 263]{brod_etal2005}.
	
	Figure \ref{fig:Figure8} shows the location, terrain surface and Registro do Araguaia 
	total-field anomaly.
	The airborne data have N-S flight lines and E-W tight lines spaced from 0.5 km and 5.0 km, 
	respectively. 
	The airborne survey covered an area of approximately $58,834$ km$^{2}$ in the Goiás state, 
	central Brazil, from the beginning of July to the end of September of 2004. 
	In the study area, there are $480,830$ observation points with vertical coordinates varying 
	from $-616.60$ m to $-358.41$ m, with average at $-426.43$ m and standard deviation $30.53$ m.
	
	To reduce the effect of over sampling along flight and tight lines, we applied a 
	block-averaging approach similar to that proposed by \citet{soler_uieda2021}. 
	We have first dividing the survey area into a regular grid of adjacent 
	horizontal blocks along the $x$ and $y$ direction, respectively. 
	To the horizontal center of each block, we assign the averaged total-field anomaly and
	vertical coordinate of the observation points within the corresponding block.
	The result is a new dataset with $N = 12,000$ averaged total-field anomaly and vertical 
	coordinates values distributed at points $\mathbf{r}_{i} = (x_{i}, y_{i}, z_{i})^{\top}$, 
	$i \in \{1, \dots, N = N_{x} \, N_{y}\}$, where $N_{x} = 150$, $N_{y} = 80$, with constant  
	horizontal spacing $dx = 297.64$ m and $dy = 475.41$ m along the $x$ and $y$ directions, 
	respectively. 
	In this case, we have an $N \times 1$ observed data vector $\mathbf{d}^{o}$ 
	(Definition \ref{def:discrete-eqlayer}) with $i$-th element $d^{o}_{i}$ given by the 
	averaged total-field anomaly at the point $\mathbf{r}_{i} = (x_{i}, y_{i}, z_{i})^{\top}$ 
	and $D_{d} = D(1, \hat{\mathbf{u}})\, D(1, \hat{\mathbf{v}})$ 
	(equations \ref{eq:observed-data-space} and \ref{eq:diff-operators-generic}), where 
	$\hat{\mathbf{u}}$ and $\hat{\mathbf{v}}$ are unit vectors defining the presumed constant 
	directions of the main field at the study area and of the total magnetization of the true 
	sources, respectively.
	
	We have set a discrete equivalent layer composed of punctual equivalent sources arranged on 
	a regular grid of $N_{x} \times N_{y}$ points 
	$\mathbf{r}''_{j} = (x''_{j}, y''_{j}, z''_{0})^{\top}$, $j \in \{1, \dots, M = N_{x} \, N_{y}\}$, 
	so that each point $\mathbf{r}''_{j}$ is located right below a point  
	$\mathbf{r}_{i}$, on the horizontal surface $\mathcal{S}''$ with constant vertical coordinate 
	$z''_{0} = 291.10$ m. This value was set as the median of averaged vertical coordinates $z_{i}$
	plus $1.5 \: dy$. 
	
	We define the elements $a_{ij}$ of matrix $\mathbf{A}$ (Definition \ref{def:discrete-eqlayer}) 
	by using the differential operator $D_{a} = D(1, \hat{\mathbf{u}})$ 
	(equation \ref{eq:diff-operators-generic}) in the space domain, where $\hat{\mathbf{u}}$ is 
	a unit vector with inclination $-18.48^{\circ}$ and declination $-18.17^{\circ}$ equal to 
	those of the main field at the study area in August 15, 2004. 
	In this case, the equivalent layer approximates the 
	$z$ component  $\frac{\partial}{\partial z}\nabla U(\mathbf{r}'')^{\top} \hat{\mathbf{v}}$
	of the magnetic induction field (Table \ref{tbl:Table2}) due to the true sources on the 
	horizontal surface $\mathcal{S}''$.
	
	Figure \ref{fig:Figure9}a shows the estimated equivalent layer $\mathbf{p}_{\epsilon}$ 
	obtained by iteratively solving a linear inverse problem in the space domain via 
	Conjugate Gradient Least Squares (CGLS) method \citep[][ p. 166]{aster_etal2019}. 
	The estimated discrete layer $\mathbf{p}_{\epsilon}$ produces a predicted potential 
	field $\mathbf{d}_{\epsilon}$ (equation \ref{eq:predicted-data-discrete-layer}) that 
	satisfactorily fit the averaged total-field anomaly in $\mathbf{d}^{o}$, as we can see 
	in the residuals shown in Figure \ref{fig:Figure9}b. The convergence curve along the 
	iterations of the CGLS method is shown in Figure \ref{fig:Figure9}d.
	
	Using the estimated layer $\mathbf{p}_{\epsilon}$, we computed the transformed potential field 
	$\mathbf{f}_{\epsilon}$ that is shown in Figure \ref{fig:Figure9}c and 
	approximates the amplitude of the magnetic 
	induction field generated by the true sources (Table~\ref{tbl:Table2}). 
	One important attribute of this field is its weak 
	dependence on the total-magnetization direction of the true sources. 
	Besides, it is computed without requiring any information about the total-magnetization
	direction of the true sources, which is a clear advantage with respect to the 
	RTP anomaly, for example. 
	The maxima of the amplitude field indicate, with considerably accuracy, the horizontal
	position of true sources.
	
	To construct $\mathbf{f}_{\epsilon}$, we first derived three intermediate transformed fields: 
	$\mathbf{f}^{x}_{\epsilon}$, $\mathbf{f}^{y}_{\epsilon}$, and $\mathbf{f}^{z}_{\epsilon}$. 
	Each vector $\mathbf{f}^{u}_{\epsilon}$ 
	($u \in {x, y, z}$) is $N \times 1$ and its $i$-th element approximates a Cartesian 
	component of the magnetic induction field (Table \ref{tbl:Table2}) at the point 
	$\mathbf{r}_{i}$.
	To compute $\mathbf{f}^{x}_{\epsilon}$, we defined an $N \times M$ transformation 
	matrix $\mathbf{T}^{x}$ (equation~\ref{eq:transformed-field-discrete-layer}) whose 
	elements are given by 
	$t^{x}_{ij} \equiv D_{t} \Psi(\mathbf{r}_{i}, \mathbf{r}''_{j})$, 
	where $D_{t} = D(1,\hat{\mathbf{x}})$ 
	(equation~\ref{eq:diff-operators-generic}), resulting in
	$t^{x}_{ij} \equiv \frac{\partial}{\partial x} \Psi(\mathbf{r}_{i}, \mathbf{r}''_{j})$. 
	Similarly, the matrices for $\mathbf{f}^{y}_{\epsilon}$ and $\mathbf{f}^{z}_{\epsilon}$ have 
	elements $t^{y}_{ij} = \frac{\partial}{\partial y} \Psi(\mathbf{r}_{i}, \mathbf{r}''_{j})$ and 
	$t^{z}_{ij} = \frac{\partial}{\partial z} \Psi(\mathbf{r}_{i}, \mathbf{r}''_{j})$, 
	respectively.
	Finally, the transformed potential field $\mathbf{f}_{\epsilon}$ was obtained by computing the 
	square root of the sum of the squares of the elements of $\mathbf{f}^{x}_{\epsilon}$, 
	$\mathbf{f}^{y}_{\epsilon}$, and $\mathbf{f}^{z}_{\epsilon}$ at each point $\mathbf{r}_{i}$.
	
	The computed $\mathbf{f}_{\epsilon}$ shown in Figure \ref{fig:Figure9}c suggests a complex 
	magnetic source. The strong red values indicate shallow portions of the source, which 
	seems to be located over a nearly circular and deeper body. 
	It seems that this complex source is located at the intersection of linear geological 
	structures having nearly NW-SE and NE-SW orientations. 
	These directions coincide with major lineaments that were mapped in continental scale and 
	represent zones of crustal weakness that were possibly reactivated in the Cretaceous by 
	the alkaline magmatism around the Paraná Basin, including that present in the GAP 
	\citep{marangoni_mantovani_2013}.
	
	% Conclusions
	\section{Conclusions}
	\label{sec:conclusions}
	
	We furnish the theoretical background for potential-field transformations
	that have been empirically made in space domain by using equivalent layers since 
	this technique was introduced in applied geophysics during the 1960s.
	Our goal here is replacing the widely accepted equivalent-layer principle, which 
	refers to the empirical generalizations
	by which most of the studies found in the geophysical literature justify the existence of
	equivalent layers allowing computing potential-field transformations.
	Specifically, we have formulated the existence of planar and continuous equivalent layers that 
	reproduce the potential-field data within a predefined error level
	as the planar equivalent-layer theorem.
	We have proven this mathematical statement by modifying the solutions of the 
	Dirichlet's and Neumann's problems on planar surfaces that may cross the sources.
	We also establish how a particular equivalent layer can be used to perform a particular 
	potential-field transformation.
	Besides, we define the theoretical link between potential-field transformations made via 
	equivalent-layer 
	technique and those directly computed from the Fourier transform of the observed 
	potential-field data. 
	In essence, we have filled the theoretical gaps in most of the existing equivalent-layer 
	methods, including 
	those presenting alternative formulations, and have also shown 
	that a planar equivalent layer is not a purely mathematical abstraction,
	but carries physical meaning and is directly related to boundary value problems of potential 
	theory.
	Results obtained from numerical simulations and also from an airborne magnetic dataset  
	strongly corroborate our theoretical findings.
	We do not present a totally new application for the equivalent-layer technique. 
	Rather, we contribute to the general understanding on why the technique has long been 
	successfully applied to perform potential-field transformations. 
	Our work opens the door for future theoretical investigations about potential-field 
	transformations via 
	equivalent-layer technique, as well as for the relationship between planar equivalent layers 
	and the true sources.	

	%\include{aknowledgements}
	
	\clearpage
	
	%\include{algorithms}
	\tabl{symbols}{Table of symbols.
}{
	\begin{center}
		\renewcommand{\arraystretch}{0.6}
		\begin{tabular}{
				>{\smaller\raggedleft\arraybackslash}m{0.18\textwidth}>{\tiny\setlength{\baselineskip}{0.5\baselineskip}\raggedright\arraybackslash}m{0.32\textwidth}
				>{\smaller\raggedleft\arraybackslash}m{0.18\textwidth}>{\tiny\setlength{\baselineskip}{0.5\baselineskip}\raggedright\arraybackslash}m{0.32\textwidth}
			}
			{\small \textbf{Symbol}} & 
			{\small \textbf{Definition}} & 
			{\small \textbf{Symbol}} & 
			{\small \textbf{Definition}} \\
			\hline
			$\mathcal{V}'_{\nu}$ & 
			Regular region: Open, connected subset of $\mathbb{R}^{3}$, with boundary defined by a closed and compact 
			surface $\partial \mathcal{V}'$ free of double points \citep[][ p. 22]{freeden_gerhards2012}. &
			$\overline{\mathcal{V}'}$ & 
			Closed region defined as the union of a finite collection of $N_{\nu}$ disjoint regular regions $\mathcal{V}'_{\nu}$ 
			and their respective boundaries $\partial\mathcal{V}'_{\nu}$, i.e., 
			$\overline{\mathcal{V}'} = \mathcal{V}' \cup \partial\mathcal{V}'$, where 
			$\mathcal{V}' = \bigcup_{\nu = 1}^{N_{\nu}} \mathcal{V}'_{\nu}$ and 
			$\partial \mathcal{V}' = \bigcup_{\nu = 1}^{N_{\nu}} \partial\mathcal{V}'_{\nu}$. 
			We consider that $\overline{\mathcal{V}'}$ contains the sources. \\
			
			$z^{t}_{0}$ and $\mathbf{r}'$ &
			Constant vertical coordinate defining the top of the sources and 
			generic point $\mathbf{r}' = (x', y', z')^{\top}$ within the sources, i.e, $\mathbf{r}' \in \overline{\mathcal{V}'}$. &
			$\sigma(\mathbf{r}')$ & 
			Piecewise continuous, integrable and bounded function that satisfies some Hölder condition \citep[][ p. 152]{kellogg1967} 
			at all points $\mathbf{r}' \in \overline{\mathcal{V}'}$ and is zero at all points $\mathbf{r} \notin \overline{\mathcal{V}'}$.
			We assume that this function represents the density (contrast) or total-magnetization intensity within the sources. \\
			
			$\mathcal{S}''$ & 
			Horizontal surface located at a constant vertical coordinate $z''_{0}$ that may cross the sources. This surface
			is used to define all equivalent layers in the present work. &
			$\mathbf{r}''$ &
			Generic point $\mathbf{r}'' = (x'', y'', z''_{0})^{\top} \in \mathcal{S}''$. \\
			
			$\mathbf{r}''_{j}$ & 
			Specific element of a finite set of $M$ points 
			$\mathbf{r}''_{j} = (x''_{j}, y''_{j}, z''_{o})^{\top} \in \mathcal{S}''$, $j \in \{ 1, \dots, M \}$. 
			These points define the locations of equivalent sources on $\mathcal{S}''$. &
			$\mathbf{r}$ &
			Generic point $\mathbf{r} = (x, y, z)^{\top}$ outside the sources ($\mathbf{r} \notin \overline{\mathcal{V}'}$). 
			It is also used to define generic observation points with potential-field data located above the sources and the surface 
			$\mathcal{S}''$, i.e., with vertical coordinates satisfying $z < z^{t}_{0}$ and $z < z''_{0}$. \\
			
			$\mathbf{r}_{i}$ & 
			Specific element of a finite set of $N$ observation points $\mathbf{r}_{i} = (x_{i}, y_{i}, z_{i})^{\top}$ 
			located above the sources and the surface $\mathcal{S}''$, i.e., $\mathbf{r}_{i} \notin \overline{\mathcal{V}'}$, 
			$i \in \{ 1, \dots, N \}$, with vertical coordinates satisfying $z_{i} < z^{t}_{0}$ and $z_{i} < z''_{0}$. 
			We consider that these points contain potential-field data.	&
			$\mathcal{S}^{\dagger}$ and $\mathcal{S}^{\ddagger}$ & 
			Horizontal surfaces located at constant vertical coordinates $z^{\dagger}_{0}$ and $z^{\ddagger}_{0}$, respectively, 
			above the sources and the surface $\mathcal{S}''$, i.e., 
			with $z^{\dagger}_{0} < z^{t}_{0}$, $z^{\dagger}_{0} < z''_{0}$, 
			$z^{\ddagger}_{0} < z^{t}_{0}$ and $z^{\ddagger}_{0} < z''_{0}$. \\
			
			$\mathbf{r}^{\dagger}$ and $\mathbf{r}^{\ddagger}$ &
			Generic points $\mathbf{r}^{\dagger} = (x^{\dagger}, y^{\dagger}, z^{\dagger}_{0})^{\top} \in \mathcal{S}^{\dagger}$ 
			and $\mathbf{r}^{\ddagger} = (x^{\ddagger}, y^{\ddagger}, z^{\ddagger}_{0})^{\top} \in \mathcal{S}^{\ddagger}$,
			both located above the sources and the surface $\mathcal{S}''$. &
			$k_{x}$, $k_{y}$, $k$ and $\mathbf{k}$ &
			The first two variables represent the wavenumbers along the $x$ and $y$ directions, respectively, 
			while $k = \sqrt{k_{x}^{2} + k_{y}^{2}}$ and $\mathbf{k} = ( \mathtt{i} \, k_{x} \, , \, \mathtt{i} \, k_{y} \, , \, k)^{\top}$
			\citep{pedersen1978}, with $\mathtt{i}$ representing the imaginary unit. \\
			
			$\mathbf{r}^{\ddagger}_{k}$ & 
			Specific element of a finite set of $L$ points 
			$\mathbf{r}^{\ddagger}_{k} = (x^{\ddagger}_{k}, y^{\ddagger}_{k}, z^{\ddagger}_{o})^{\top} \in \mathcal{S}^{\ddagger}$, 
			$k \in \{ 1, \dots, L \}$. &
			$U$ and $U_{\varepsilon}$ & 
			Newtonian potential produced by the true sources at points $\mathbf{r} \notin \overline{\mathcal{V}'}$ 
			(equation \ref{eq:potential-U-space}) and smooth approximation for $U$ at points located above the sources
			(equation \ref{eq:potential-U-epsilon-space}). \\
			
			$\mathcal{F} [ \cdot ]$ &
			Two-dimensional Fourier transform of a generic scalar function computed on a horizontal plane, with respect to
			the spatial coordinates $x$ and $y$. To simplify the notation, we conveniently omit the dependence of 
			$\mathcal{F} [ \cdot ]$ on the wavenumbers $k_{x}$ and $k_{y}$. &
			$D_{g}$ and $\alpha_{g}$ &
			Space- and Fourier-domain representations of the differential operator acting on a given function to produce
			a generic function $g$. For example, $g(\cdot) = D_{g} U(\cdot)$ in the space domain and 
			$\mathcal{F} [ g(\cdot) ] = \alpha_{g} \mathcal{F} [ U(\cdot) ]$ in the Fourier domain. 
			$D_{g}$ and $\alpha_{g}$ are usually defined in terms of $D$ and $\alpha$ (equation \ref{eq:diff-operators-generic}). \\
			
			$k_{\varepsilon}$ and $\mu_{\varepsilon}$ &
			Fixed truncating wavenumber and fixed positive regularization parameter defining the approximating potential 
			$U_{\varepsilon}$ according to equations \ref{eq:potential-U-epsilon-Fourier-dagger-1} and 
			\ref{eq:potential-U-epsilon-Fourier-dagger-2}, respectively. &
			$\hat{\mathbf{u}}$, $\hat{\mathbf{v}}$ and $\hat{\mathbf{h}}$ &
			Constant unit vectors used in the present work.			
		\end{tabular}
	\end{center}
	\label{tbl:symbols}
}
	\renewcommand{\figdir}{figures} % figure directory

\plot{sources}{width=0.8\textwidth}{
	{2D schematic representation of sources located in subsurface, over the
	domain $\overline{\mathcal{V}} = \bigcup_{\nu = 1}^{N_{\nu}}$, $N_{\nu} = 2$, 
	with top limited by the plane $z^{t}_{0}$. 
	Points within the sources are represented by $\mathbf{r}' = (x', y', z')^{\top}$.
	The dashed line represents the surface containing the observation points 
	$\mathbf{r} = (x, y, z)^{\top}$. The thick line represents the Earth's surface.
	The planar surface $\mathcal{S}^{\dagger}$ containing points 
	$\mathbf{r}^{\dagger} = (x^{\dagger}, y^{\dagger}, z^{\dagger}_{0})^{\top}$ 
	above the sources is represented by a continuous straight line.
	The fine dashed lines represent two possible definitions of a horizontal surface $\mathcal{S}''$ 
	at a constant vertical coordinate $z''_{0}$ satisfying $z''_{0} > z^{t}_{0}$ or $z''_{0} < z^{t}_{0}$. 
	}
	\label{fig:geometry}
}

\plot{model_view_1.png}{width=0.8\textwidth}{
	{Simulated magnetic source formed by horizontally juxtaposed rectangular prisms.}
	\label{fig:model}
}

\plot{discrete_layers.png}{width=0.8\textwidth}{
	{
		Discrete equivalent layers estimated by running Algorithm \ref{alg:numerical-eqlayer}.
		(a) Computed layer having the constant magnetization direction equal to that of the 
		model shown in Figure \ref{fig:model}.
		(b) Computed layer having the constant magnetization direction different from that of the 
		model shown in Figure \ref{fig:model}.
		All values are in $\mu$T m.	
		The black line represents the projection of the model on the $xy$ plane.	
	}
	\label{fig:discrete-layers}
}

\plot{Neumann_bx_by_bz.png}{width=0.8\textwidth}{
	{
		Potential fields produced by the model shown in Figure \ref{fig:model} and by the discrete 
		layer shown in Figure \ref{fig:discrete-layers}a.
		(a), (d) and (g) show the potential fields produced by the discrete layer.
		(b), (e) and (h) show the potential fields produced by the model.
		(c), (f) and (i) show the difference between the fields produced by the layer and the model.
		The upper, middle and bottom rows show the $x$-, $y$- and $z$-components of the 
		magnetic induction field, respectively. All values are in nT.
		The black line represents the projection of the model on the $xy$ plane.	
	}
	\label{fig:Neumann-bx-by-bz}
}

\plot{Neumann_atfa_b_rtp.png}{width=0.8\textwidth}{
	{
		Potential fields produced by the model shown in Figure \ref{fig:model} and by the discrete 
		layer shown in Figure \ref{fig:discrete-layers}a.
		(a), (d) and (g) show the potential fields produced by the discrete layer.
		(b), (e) and (h) show the potential fields produced by the model.
		(c), (f) and (i) show the difference between the fields produced by the layer and the model.
		The upper, middle and bottom rows show the approximated total-field anomaly, amplitude of 
		magnetic induction field and RTP anomaly, respectively. All values are in nT.
		The black line represents the projection of the model on the $xy$ plane.		
	}
	\label{fig:Neumann-atfa-b-rtp}
}

\plot{generic_bx_by_bz.png}{width=0.8\textwidth}{
	{
		Potential fields produced by the model shown in Figure \ref{fig:model} and by the discrete 
		layer shown in Figure \ref{fig:discrete-layers}b.
		(a), (d) and (g) show the potential fields produced by the discrete layer.
		(b), (e) and (h) show the potential fields produced by the model.
		(c), (f) and (i) show the difference between the fields produced by the layer and the model.
		The upper, middle and bottom rows show the $x$-, $y$- and $z$-components of the 
		magnetic induction field, respectively. All values are in nT.
		The black line represents the projection of the model on the $xy$ plane.	
	}
	\label{fig:dipole-bx-by-bz}
}

\plot{generic_atfa_b_rtp.png}{width=0.8\textwidth}{
	{
		Potential fields produced by the model shown in Figure \ref{fig:model} and by the discrete 
		layer shown in Figure \ref{fig:discrete-layers}b.
		(a), (d) and (g) show the potential fields produced by the discrete layer.
		(b), (e) and (h) show the potential fields produced by the model.
		(c), (f) and (i) show the difference between the fields produced by the layer and the model.
		The upper, middle and bottom rows show the approximated total-field anomaly, amplitude of 
		magnetic induction field and RTP anomaly, respectively. All values are in nT.
		The black line represents the projection of the model on the $xy$ plane.		
	}
	\label{fig:dipole-atfa-b-rtp}
}

\plot{generic_rtp_fixed.png}{width=0.8\textwidth}{
	{
		RTP anomalies produced by the model shown in Figure \ref{fig:model} and by the discrete 
		layer shown in Figure \ref{fig:discrete-layers}b, after correction.
		(a) shows the fixed RTP anomaly produced by the discrete layer.
		(b) shows the RTP anomaly produced by the model.
		(c) shows the difference between the RTP anomalies produced by the layer and the model.
		All values are in nT.
		The black line represents the projection of the model on the $xy$ plane.		
	}
	\label{fig:dipole-rtp-fixed}
}


	
	\clearpage
	
	% Appendices
	\appendix
	
	% Appendix A
	\append{Some useful upper bounds}
	\label{app:upper-bounds}
	
	Here we define upper bounds for the absolute values
	of $\mathcal{F} [ U(\mathbf{r}^{\dagger}) ]$ and $\mathcal{F} [ d^{o}(\mathbf{r}^{\dagger}) ]$ 
	(equations \ref{eq:potential-U-Fourier} and \ref{eq:observed-data-Fourier}), 
	i.e., for the amplitude spectra of $U(\mathbf{r})$ and $d^{o}(\mathbf{r})$ 
	(equations \ref{eq:potential-U-space} and \ref{eq:observed-data-space}). 
	To do this, we have adapted the approach presented by
	\citet{bhattacharyya1967} to define an upper bound for the amplitude spectrum of 
	$\tfrac{\partial}{\partial z} U(\mathbf{r})$. 
	We first stress that 
	$e^{- k \, (z' - z^{\dagger}_{0})} \le e^{- k \, (z^{t}_{0} - z^{\dagger}_{0})}$ because
	$z^{t}_{0} \le z'$, where $z^{t}_{0}$ defines the top of the sources and $z'$ the vertical
	coordinate of a point inside the sources. Using this identity in  
	$\mathcal{F} [ U(\mathbf{r}^{\dagger}) ]$ (equation \ref{eq:potential-U-Fourier}), we obtain
	\begin{equation}
		\vert \mathcal{F} [ U(\mathbf{r}^{\dagger}) ] \vert \le 2 \pi \kappa 
		\frac{e^{- k \, \Delta z}}{k}
		\int\limits_{z^{t}_{0}}^{\infty}
		\mathcal{F} [ \sigma(\mathbf{r}') ] \:
		\: dz' \: , \quad \Delta z = z^{t}_{0} - z^{\dagger}_{0} > 0 \: .
		\label{eq:potential-U-spectrum-bound1}
	\end{equation}
	This equation is valid because we assume that $\sigma(\mathbf{r}')$ (Table \ref{tbl:Table1}) 
	is an integrable scalar function that is zero at all points $\mathbf{r}' \notin \overline{\mathcal{V}}$, 
	where $\overline{\mathcal{V}}$ has a finite extension. 
	Manipulating $\mathcal{F} [ \sigma(\mathbf{r}') ]$ (equation \ref{eq:sigma-Fourier}), 
	we obtain
	\begin{equation}
		\vert \mathcal{F} [ \sigma(\mathbf{r}') ] \vert \le
		\int\limits_{-\infty}^{\infty}\int\limits_{-\infty}^{\infty}
		\vert \sigma(\mathbf{r}') \vert 
		\: dx'dy'
		\label{eq:sigma-spectrum-bound}
	\end{equation}
	and, consequently, that
	\begin{equation}
		2 \pi \kappa \, 
		\int\limits_{z^{t}_{0}}^{\infty}
		\vert \mathcal{F} [ \sigma(\mathbf{r}') ] \vert
		\: dz' \le \Omega \: ,
		\label{eq:Omega-bound}
	\end{equation}
	where
	\begin{equation}
		\Omega = 2 \pi \kappa \, \iiint\limits_{\mathcal{V}}
		\vert \sigma(\mathbf{r}') \vert 
		\: d\mathcal{V}'
		\label{eq:Omega}
	\end{equation}
	is a constant depending only on the position and shape of the sources. 
	Then, by combining equations 
	\ref{eq:potential-U-spectrum-bound1}--\ref{eq:Omega},
	we conclude that the amplitude spectrum of 
	$U(\mathbf{r})$ (equation \ref{eq:potential-U-space}) on a horizontal plane 
	$z^{\dagger}_{0} < z^{t}_{0}$ located above the sources satisfies 
	\begin{equation}
		\vert \mathcal{F} [ U(\mathbf{r}^{\dagger}) ] \vert \le
		\frac{e^{- k \, \Delta z}}{k} \, \Omega \: ,
		\: \quad \Delta z = z^{t}_{0} - z^{\dagger}_{0} > 0 \: .
		\label{eq:potential-U-spectrum-bound}
	\end{equation}
	From the upper bound of $\vert \mathcal{F} [ U(\mathbf{r}^{\dagger}) ] \vert$
	(equation \ref{eq:potential-U-spectrum-bound}) and $\mathcal{F} [ d^{o}(\mathbf{r}^{\dagger}) ]$
	(equation \ref{eq:observed-data-Fourier}), we also deduce
	\begin{equation}
		\vert \mathcal{F} [ d^{o}(\mathbf{r}^{\dagger}) ] \vert \le
		\vert \alpha_{d} \vert \, \Omega \, 
		\frac{e^{- k \, \Delta z}}{k} \: ,
		\: \quad \Delta z = z^{t}_{0} - z^{\dagger}_{0} > 0 \: .
		\label{eq:observed-data-spectrum-bound}
	\end{equation}
	Note that, for the case in which $\alpha_{d} = k$ (equation \ref{eq:observed-data-Fourier}), 
	we obtain the same result presented by \citet{bhattacharyya1967} as the upper bound for the 
	amplitude spectrum of $\tfrac{\partial}{\partial z} U(\mathbf{r})$, whereas
	$\alpha_{d} = 1$ leads to the upper bound defined by equation \ref{eq:potential-U-spectrum-bound}.
	Equations \ref{eq:potential-U-spectrum-bound} and \ref{eq:observed-data-spectrum-bound} 
	show that the amplitude spectra 
	$\vert \mathcal{F} [ U(\mathbf{r}^{\dagger}) ] \vert$ and 
	$\vert \mathcal{F} [ d^{o}(\mathbf{r}^{\dagger}) ] \vert$ on a plane $z^{\dagger}_{0} < z^{t}_{0}$ 
	above the sources 
	(i) are finite at all wavenumbers $k = \sqrt{k_{x}^{2} + k_{y}^{2}}$ and (ii) tend
	to zero as $k$ goes to infinite.
	These results confirm the well-known fact that 
	$\mathcal{F} [ U(\mathbf{r}^{\dagger}) ]$ and $\mathcal{F} [ d^{o}(\mathbf{r}^{\dagger}) ]$ 
	(equations \ref{eq:potential-U-Fourier} and \ref{eq:observed-data-Fourier}) are square integrable.
	
	% Appendix B
	\append{Phase spectrum of the potentials $U$ and $U_{\varepsilon}$}
	\label{app:potential-U-min-phase}
	
	Let the phase spectrum \cite[e.g.,][ p. 261]{blakely1996} of  
	$\mathcal{F} [ U(\mathbf{r}^{\dagger}) ]$ (equation \ref{eq:potential-U-Fourier}) 
	on a horizontal plane $z^{\dagger}_{0}$ above the sources be defined as follows:
	\begin{equation}
		\Theta [ U(\mathbf{r}^{\dagger}) ] = \tan^{-1} \left( 
		\frac{\text{Im} \, \mathcal{F} [ U(\mathbf{r}^{\dagger}) ]}{\text{Re} \, \mathcal{F} [ U(\mathbf{r}^{\dagger}) ]} 
		\right) \: ,
		\label{eq:potential-U-phase-spectrum}
	\end{equation}
	where ``$\text{Re}$'' and ``$\text{Im}$'' denote the real and imaginary parts, respectively.
	
	According to equation \ref{eq:potential-U-phase-spectrum}, the phase spectrum of $U(\mathbf{r}^{\dagger})$ (equation \ref{eq:potential-U-Fourier}) 
	represents a weighted sum of the Fourier transforms $\mathcal{F} [ \sigma(\mathbf{r}') ]$
	(equation \ref{eq:sigma-Fourier}) over all depth slices crossing the sources,
	where the weights are real functions that exponentially decay with depth $z'$.
	It means that $\Theta [ U(\mathbf{r}^{\dagger}) ]$ attributes more weight to the shallow than to the deeper parts of the sources. 
	Because of that, the phase spectrum of $U(\mathbf{r}^{\dagger})$ approximates that of $\sigma(\mathbf{r}')$ and, consequently, 
	$U(\mathbf{r}^{\dagger})$ has approximately the same horizontal spatial distribution of the (shallow parts of) true sources.
	The same reasoning can be applied to the phase spectra of any vertical integral or vertical derivative of the potential $U$.
	
	It can be easily verified that, for the case in which $U_{\varepsilon}$ is given by equation 
	\ref{eq:potential-U-epsilon-Fourier-dagger-1} in the Fourier domain, we have that
	\begin{equation}
		\Theta [ U_{\varepsilon}(\mathbf{r}'') ] = \begin{cases}
			\Theta [ U(\mathbf{r}^{\dagger}) ] & , \quad 0 < k \le k_{\varepsilon} \\
			0 & , \quad k > k_{\varepsilon}
		\end{cases} \: .
		\label{eq:potential-U-epsilon-1-phase-spectrum}
	\end{equation}
	It can also be easily verified that, for the case in which $U_{\varepsilon}$ is given by equation 
	\ref{eq:potential-U-epsilon-Fourier-dagger-2} in the Fourier domain, we have that
	\begin{equation}
		\Theta [ U_{\varepsilon}(\mathbf{r}'') ] = \Theta [ U(\mathbf{r}^{\dagger}) ] \: .
		\label{eq:potential-U-epsilon-2-phase-spectrum}
	\end{equation}
	
	\clearpage
	
	% References
	\bibliographystyle{seg}  % style file is seg.bst
	\bibliography{references}
	
	\clearpage
	
\end{document}
